\chapter{The Quot scheme}
\label{chap:Picard_QuotScheme}
% archon:covers AlgebraicJacobian/Picard/QuotScheme.lean

% STRATEGY NOTE (iter-174). This chapter is the A.2.b sub-row of Route~A
% (positive-genus arm of nonempty_jacobianWitness). It packages the
% Grothendieck--Altman--Kleiman Quot-scheme construction
% \(\mathrm{Quot}^{\Phi,L}_{E/X/S}\) -- a projective \(S\)-scheme representing the
% functor of \(T\)-flat coherent quotients of \(E_T\) on \(X_T = X \times_S T\) with
% Hilbert polynomial \(\Phi\) on every fiber. The Quot construction is the
% \emph{existence engine} behind the relative Picard functor: A.2.c
% (\cref{chap:Picard_FGAPicRepresentability}) wires this chapter together with
% the relative Picard functor (\cref{chap:Picard_RelPicFunctor}, A.1.c)
% through the Abel map \(\mathrm{Div}_{C/k} \to \Pic^\sharp_{(C/k)\et}\), which factors
% through the Hilbert scheme \(\Hilb_{C/k} = \Quot_{\mathcal{O}_C/C/k}\) -- the
% case \(E = \mathcal{O}_X\) of the Quot functor. This chapter is gated on
% A.2.a (\cref{thm:flattening_stratification_exists}, written this iter) and
% A.1.a (\cref{chap:Picard_RelativeSpec}, on disk).
%
% LOAD-BEARING SUB-PREREQUISITE. The construction embeds Quot as a
% locally closed subfunctor of a Grassmannian
% \(\mathrm{Grass}(W \otimes_{\mathcal{O}_S} \mathrm{Sym}^r V, \Phi(r))\).
% Mathlib (at the pinned commit) carries no Grassmannian \emph{scheme} -- only
% a linear-algebra Grassmannian as a finite-rank-quotient variety. The
% scheme-level Grassmannian over an arbitrary noetherian base is therefore an
% in-project sub-build. This chapter introduces it as
% \cref{def:grassmannian_scheme} and states its representability; it is a
% candidate for promotion to its own chapter
% \texttt{Picard\_GrassmannianScheme.tex} (writer iter-175+) -- see "Notes
% for Plan Agent" in the writer report.

\section{Setup and motivation}
\label{sec:quot_setup}

Throughout this chapter, let \(S\) be a noetherian scheme, \(\pi : X \to S\) a
projective morphism, \(L\) a relatively very ample line bundle on \(X\) over \(S\),
and \(E\) a coherent sheaf on \(X\). The Quot construction packages the
\(T\)-flat coherent quotients \(E_T \twoheadrightarrow \mathcal{F}\) of the
pull-back \(E_T\) to \(X_T = X \times_S T\), modulo the equivalence
\(\langle \mathcal{F}, q\rangle = \langle \mathcal{F}', q'\rangle\) defined by
\(\ker(q) = \ker(q')\). Fixing the Hilbert polynomial
\(\Phi(\lambda) = \chi(\mathcal{F}|_{X_t}(\lambda)) \in \mathbb{Q}[\lambda]\)
(constant on connected \(T\) by flatness)
stratifies the Quot functor into pieces \(\Quot^{\Phi,L}_{E/X/S}\), each of
which Grothendieck proved is representable by a projective \(S\)-scheme.

% SOURCE: [Nitsure], "Construction of Hilbert and Quot Schemes", \S 1 and
% \S 5; FGA Explained Ch.~5, AMS MSM 123 (read from
% references/nitsure-hilbert-quot-src/nitsure-hilbert-quot.tex, L287-L500 +
% L2154-L2542). The Hartshorne reference for the polynomial-eventually
% property of Euler characteristics is [Hartshorne], III.5.2 (the scanned
% PDF at references/hartshorne-algebraic-geometry.pdf has no text layer, so
% the verbatim source-quote in \ is taken from
% Nitsure \S 1; Hartshorne III.5.2 is cited in prose as the textbook
% reference).

\section{Hilbert polynomial of a coherent sheaf}
\label{sec:hilbert_polynomial}

\begin{definition}\leanok
[Hilbert polynomial]
  \label{def:hilbert_polynomial}
  \lean{AlgebraicGeometry.Scheme.hilbertPolynomial}
  % SOURCE: [Nitsure], \S 1, sub-section "Stratification by Hilbert
  % Polynomials" (read from
  % references/nitsure-hilbert-quot-src/nitsure-hilbert-quot.tex, L453-L478).
  % SOURCE QUOTE:
  % "Let \(X\) be a finite type scheme over
  %  a field \(k\), together with a line bundle \(L\).
  %  Recall that if \(F\) is a coherent sheaf on
  %  \(X\) whose support is proper over \(k\), then
  %  the \textbf{Hilbert polynomial} \({\Phi}\in \Q[\lambda]\) of \(F\) is defined by
  %  the function
  %  \(\){\Phi}(m) = \chi(F(m)) =
  %  \sum_{i=0}^n (-1)^i \dim_k H^i(X, F\otimes L^{\otimes m})\(\)
  %  where the dimensions of the cohomologies are finite because of the
  %  coherence and properness conditions. The fact that \(\chi(F(m))\)
  %  is indeed a polynomial in \(m\) under the above assumption is a special
  %  case of what is known as Snapper's Lemma (see Kleiman [K] for a proof).
  %
  %  Let \(X\to S\) be a finite type morphism of noetherian schemes,
  %  and let \(L\) be a line bundle on \(X\). Let \(\F\)
  %  be any coherent sheaf on
  %  \(X\) whose schematic support is proper over \(S\). Then for each \(s\in S\),
  %  we get a polynomial \({\Phi}_s \in  \Q[\lambda]\)
  %  which is the Hilbert polynomial of
  %  the restriction \(\F_s = F|_{X_s}\) of \(\F\) to the fiber \(X_s\) over \(s\),
  %  calculated with respect to the line bundle \(L_s = L|_{X_s}\).
  %  If \(\F\) is flat over \(S\) then the function \(s\mapsto {\Phi}_s\)
  %  from the set of points of \(S\) to the polynomial ring \(\Q[\lambda]\) is
  %  known to be locally constant on \(S\)."
  \textit{Source: [Nitsure], \S 1 (FGA Explained Ch.~5); cf.\ [Hartshorne],
  III.5.2.}
  Let \(X \to S\) be a finite-type morphism of noetherian schemes and let \(L\)
  be a line bundle on \(X\). For a coherent sheaf \(\mathcal{F}\) on \(X\) whose
  schematic support is proper over \(S\) and a point \(s \in S\), the
  \emph{Hilbert polynomial of \(\mathcal{F}\) at \(s\) relative to \(L\)} is the
  polynomial \(\Phi_{\mathcal{F},s} \in \mathbb{Q}[\lambda]\) characterised by
  \[
    \Phi_{\mathcal{F},s}(m)
    \;=\; \chi(X_s,\, \mathcal{F}|_{X_s} \otimes L_s^{\otimes m})
    \;=\; \sum_{i \ge 0} (-1)^i \dim_{\kappa(s)} H^i(X_s,\, \mathcal{F}|_{X_s} \otimes L_s^{\otimes m})
  \]
  (a polynomial in \(m\) by Snapper's Lemma; the sum is finite by coherence and
  properness). When \(\mathcal{F}\) is \(S\)-flat the function
  \(s \mapsto \Phi_{\mathcal{F},s}\) is locally constant on \(S\), so \(\Phi\) is
  well-defined on connected components.
\end{definition}

\section{The Quot functor}
\label{sec:quot_functor}

\begin{definition}\leanok
[Quot functor]
  \label{def:quot_functor}
  \lean{AlgebraicGeometry.Scheme.QuotFunctor}
  \uses{def:hilbert_polynomial}

  % SOURCE: [Nitsure], \S 1, sub-section "The Functors \(\Hilb_{X/S}\) and
  % \(\Quot_{E/X/S}\)" (read from
  % references/nitsure-hilbert-quot-src/nitsure-hilbert-quot.tex, L394-L433
  % + L481-L498).
  % SOURCE QUOTE:
  % "Let \(S\) be a noetherian scheme and let \(X\to S\) be a finite type
  %  scheme over it. Let \(E\) be a coherent sheaf on \(X\).
  %  Let \(Sch_S\) denote the category of all locally noetherian schemes
  %  over \(S\). For any \(T\to S\) in \(Sch_S\), a \textbf{family of quotients of \(E\)
  %  parametrised by \(T\)} will mean
  %  a pair \((\F,q)\) consisting of
  %
  %  (i) a coherent sheaf \(\F\) on \(X_T = X\times_ST\)
  %  such that the schematic support of \(\F\) is proper over \(T\)
  %  and \(\F\) is flat over \(T\), together with
  %
  %  (ii) a surjective \({\mathcal O}_{X_T}\)-linear homomorphism of sheaves
  %  \(q:E_T  \to \F\) where \(E_T\) is the pull-back of \(E\) under
  %  the projection \(X_T \to X\).
  %
  %  Two such families \((\F,q)\) and \((\F,q)\) parametrised by \(T\) will
  %  be regarded as equivalent if \(\ker(q) = \ker(q')\)),
  %  and \(\langle \F, q \rangle\) will denote
  %  an equivalence class. Then as properness and flatness are preserved
  %  by base-change, and as tensor-product is right exact,
  %  the pull-back of \(\langle \F, q \rangle\) under an \(S\)-morphism \(T'\to T\) is
  %  well-defined, which gives
  %  a set-valued contravariant functor
  %  \(\Quot_{E/X/S} : Sch_S \to Sets\) under which
  %  \(\)T\mapsto  \{\mbox{ All } \langle \F, q \rangle \mbox{ parametrised by }T\}\(\)
  %  [...]
  %  This shows that the functor \(\Quot_{E/X/S}\) naturally decomposes
  %  as a co-product
  %  \(\)\Quot_{E/X/S} = \coprod_{{\Phi}\in \Q[\lambda]}\, \Quot_{E/X/S}^{\Phi,L}\(\)
  %  where for any polynomial \({\Phi}\in \Q[\lambda]\), the functor
  %  \(\Quot_{E/X/S}^{\Phi,L}\) associates
  %  to any \(T\) the set of all equivalence classes of
  %  families \(\langle \F,q \rangle\) such that at each \(t\in T\) the
  %  Hilbert polynomial of the restriction \(\F_t\), calculated
  %  using the pull-back of \(L\), is \({\Phi}\)."
  \textit{Source: [Nitsure], \S 1 (FGA Explained Ch.~5).}
  Let \(S\) be a noetherian scheme, \(\pi : X \to S\) a finite-type morphism,
  \(L\) a line bundle on \(X\), \(E\) a coherent sheaf on \(X\), and
  \(\Phi \in \mathbb{Q}[\lambda]\). The \emph{Quot functor with Hilbert
  polynomial \(\Phi\)}
  \[
    \Quot^{\Phi,L}_{E/X/S} : (\mathrm{Sch}/S)^{op} \longrightarrow \mathbf{Set}
  \]
  sends an \(S\)-scheme \(T \to S\) to the set of equivalence classes
  \(\langle \mathcal{F}, q\rangle\) of pairs \((\mathcal{F}, q)\) where
  \begin{enumerate}
    \item \(\mathcal{F}\) is a coherent sheaf on \(X_T = X \times_S T\) whose
      schematic support is proper over \(T\) and which is flat over \(T\);
    \item \(q : E_T \twoheadrightarrow \mathcal{F}\) is a surjective
      \(\mathcal{O}_{X_T}\)-linear homomorphism, where \(E_T\) is the pull-back
      of \(E\) along \(X_T \to X\);
    \item for every \(t \in T\) the Hilbert polynomial of \(\mathcal{F}|_{X_t}\)
      with respect to \(L|_{X_t}\) (\cref{def:hilbert_polynomial}) equals
      \(\Phi\).
  \end{enumerate}
  Two pairs \((\mathcal{F}, q)\) and \((\mathcal{F}', q')\) are equivalent iff
  \(\ker(q) = \ker(q')\). On morphisms \(f : T' \to T\) over \(S\), the functor
  acts by pulling back the quotient along \(X_{T'} \to X_T\) (well-defined
  since tensor product is right-exact and preserves both flatness and
  properness). The Hilbert scheme is the special case
  \(\Hilb^{\Phi,L}_{X/S} := \Quot^{\Phi,L}_{\mathcal{O}_X/X/S}\).
\end{definition}

\section{The Grassmannian scheme}
\label{sec:grassmannian_scheme}

The Quot construction proceeds by embedding \(\Quot^{\Phi,L}_{E/X/S}\) as a
locally closed subfunctor of a Grassmannian over \(S\). Since Mathlib carries
no Grassmannian \emph{scheme} (only a finite-rank-quotient variety in
linear algebra), we package the construction here.

\begin{definition}\leanok
[Grassmannian scheme]
  \label{def:grassmannian_scheme}
  \lean{AlgebraicGeometry.Scheme.Grassmannian}
  \uses{def:quot_functor}
  % SOURCE: [Nitsure], \S 1, "Elementary Examples, Exercises~(2)" and
  % "Construction of Grassmannian" (read from
  % references/nitsure-hilbert-quot-src/nitsure-hilbert-quot.tex, L546-L599
  % + L773-L850).
  % SOURCE QUOTE:
  % "For any integers \(r\ge d\ge 1\), an explicit construction
  %  the Grassmannian scheme
  %  \(\grass(r,d)\) over \(\Z\), together with the tautological
  %  quotient \(u:\oplus^r\,{\mathcal O}_{\grass(r,d)} \to \U\) where \
  %  \(\U\) is a rank \(d\) locally free sheaf on \(\grass(r,d)\),
  %  has been given at the end of this section. [...]
  %  Show that \(\grass(r,d)\) together with the quotient
  %  \(u:\oplus^r\,{\mathcal O}_{\grass(r,d)} \to \U\) represents the
  %  contravariant functor
  %  \(\)\Grass(r,d) = \Quot_{\oplus^r\,{\mathcal O}_{\Z}/\Z/\Z}^{d, {\mathcal O}_{\Z}}\(\)
  %  from schemes to sets, which associates to any \(T\) the
  %  set of all equivalence classes
  %  \(\langle \F, q\rangle\) of quotients \(q: \oplus^r\,{\mathcal O}_T \to \F\)
  %  where \(\F\) is a locally free sheaf on \(T\) of rank \(d\).
  %  [...]
  %  Using the above show that, more generally, if \(S\) is a scheme and
  %  \(E\) is a locally free \({\mathcal O}_S\)-module of rank \(r\),
  %  the functor \(\Grass(E,d)=\Quot_{E/S/S}^{d,{\mathcal O}_S}\)
  %  on all \(S\)-schemes which
  %  by definition associates to any \(T\)
  %  the set of all equivalence classes \(\langle \F, q\rangle\) of quotients
  %  \(q: E_T \to \F\) where \(\F\) is a locally free sheaf on \(T\) of rank \(d\),
  %  is representable. The representing scheme
  %  is denoted by \(\grass(E,d)\) and is called the rank \(d\)
  %  relative Grassmannian of \(E\) over \(S\)."
  \textit{Source: [Nitsure], \S 1, Exercise (2) and "Construction of
  Grassmannian" (FGA Explained Ch.~5).}
  Let \(S\) be a noetherian scheme and let \(V\) be a locally free
  \(\mathcal{O}_S\)-module of rank \(r\). For an integer \(1 \le d \le r\), the
  \emph{Grassmannian of rank-\(d\) quotients of \(V\)} is the functor
  \[
    \mathrm{Grass}(V, d) : (\mathrm{Sch}/S)^{op} \longrightarrow \mathbf{Set},
    \qquad
    T \longmapsto \{ \langle \mathcal{F}, q\rangle :
    q : V_T \twoheadrightarrow \mathcal{F},\; \mathcal{F}
    \text{ locally free of rank } d \}
  \]
  with equivalence \(\ker(q) = \ker(q')\). Concretely \(\mathrm{Grass}(V, d) =
  \Quot^{d,\mathcal{O}_S}_{V/S/S}\) (the Quot functor for \(X = S\), \(E = V\),
  constant Hilbert polynomial \(\Phi = d\)).
\end{definition}

\begin{theorem}\leanok
[Representability of the Grassmannian]
  \label{thm:grassmannian_representable}
  \lean{AlgebraicGeometry.Scheme.Grassmannian.representable}
  \uses{def:grassmannian_scheme, thm:relative_spec_exists, thm:relative_spec_univ}

  % SOURCE: [Nitsure], \S 1, sub-section "Construction by gluing together
  % affine patches" through "Projective embedding"
  % (read from references/nitsure-hilbert-quot-src/nitsure-hilbert-quot.tex,
  % L798-L930).
  % SOURCE QUOTE:
  % "For any integers \(r\ge d\ge 1\), the Grassmannian scheme
  %  \(\grass(r,d)\) over \(\Z\), together with the tautological
  %  quotient \(u:\oplus^r\,{\mathcal O}_{\grass(r,d)} \to \U\) where \
  %  \(\U\) is a rank \(d\) locally free sheaf on \(\grass(r,d)\),
  %  can be explicitly constructed as follows. [...]
  %  the schemes
  %  \(U^I\), as \(I\) varies over all the
  %  \({r\choose d}\) different
  %  subsets of  \(\{ 1,\ldots, r\}\) of cardinality \(d\),
  %  can be glued together by the co-cycle \((\theta_{I,J})\) to form a
  %  finite-type scheme \(\grass(r,d)\) over \(\Z\). As each \(U^I\) is
  %  isomorphic to \(\AA^{d(r-d)}_{\Z}\), it follows that
  %  \(\grass(r,d) \to \spec \Z\) is smooth of relative dimension \(d(r-d)\).
  %  [...]
  %  We now show that \(\pi : \grass(r,d) \to \spec \Z\) is proper. It is enough to
  %  verify the valuative criterion of properness for discrete valuation rings.
  %  [...]
  %  As \(\U\) is given by the transition functions \(g_{I,J}\) described above,
  %  the determinant line bundle
  %  \(\det(\U)\) is given by the transition functions
  %  \(\det(g_{I,J}) = 1/P^I_J \in GL_1(U^I_J)\).
  %  [...]
  %  We leave it to the reader to verify that the sections \(\sigma_I\) form
  %  a linear system which is base point free and separates points
  %  relative to \(\spec \Z\), and so gives an embedding of
  %  \(\grass(r,d)\) into \(\PP^m_{\Z}\) where \(m = {r\choose d}-1\).
  %  This is a closed embedding by the properness of
  %  \(\pi:\grass(r,d) \to \spec \Z\). In particular, \(\det(\U)\)
  %  is a relatively very ample line bundle on \(\grass(r,d)\) over \(\Z\)."
  \textit{Source: [Nitsure], \S 1, "Construction of Grassmannian" (FGA
  Explained Ch.~5).}
  Let \(S\) be a noetherian scheme, \(V\) a locally free \(\mathcal{O}_S\)-module
  of rank \(r\), and \(1 \le d \le r\). The functor \(\mathrm{Grass}(V,d)\) of
  \cref{def:grassmannian_scheme} is representable by a smooth projective
  \(S\)-scheme \(\mathrm{Gr}_S(V, d) \to S\) of relative dimension \(d(r-d)\),
  equipped with a tautological quotient
  \(\pi^* V \twoheadrightarrow \mathcal{U}\) of rank \(d\). The determinant
  line bundle \(\det(\mathcal{U})\) is relatively very ample, and the
  associated linear system gives a closed embedding
  \(\mathrm{Gr}_S(V, d) \hookrightarrow
  \mathbb{P}_S\!\left(\bigwedge^d V\right)\) -- the \emph{Pl\"ucker embedding}.
\end{theorem}

% SOURCE QUOTE PROOF: see the "Construction of Grassmannian" sub-section of
% [Nitsure], \S 1 (references/nitsure-hilbert-quot-src/nitsure-hilbert-quot.tex,
% L773-L930), which gives the construction in five named steps
% ("Construction by gluing together affine patches", "Separatedness",
% "Properness", "Universal quotient", "Projective embedding"). The total
% verbatim text is ~150 lines; it is reproduced piecewise in the
% SOURCE QUOTE comment of the theorem statement above.
\begin{proof}
  
  Over \(S = \Spec \mathbb{Z}\) and \(V = \mathcal{O}_S^{\oplus r}\) the
  Grassmannian \(\mathrm{Gr}(r, d)\) is constructed by gluing
  \(\binom{r}{d}\) affine charts \(U^I = \Spec\mathbb{Z}[X^I]\) indexed by
  size-\(d\) subsets \(I \subseteq \{1, \dots, r\}\). The chart \(U^I\)
  carries the matrix \(X^I\) of size \(d \times r\) whose \(I\)-th \(d \times d\)
  minor is the identity, with the remaining \(d(r-d)\) entries free
  variables (so \(U^I \cong \mathbb{A}^{d(r-d)}_{\mathbb{Z}}\)). The
  transition maps \(\theta_{I,J} : U^I \cap U^J \to U^J \cap U^I\) defined
  by \(X^J \mapsto (X^I_J)^{-1} X^I\) on the locus where \(\det X^I_J\) is
  invertible satisfy the cocycle condition
  \(\theta_{I,K} = \theta_{I,J} \circ \theta_{J,K}\). The glued
  \(\mathrm{Gr}(r,d) \to \Spec\mathbb{Z}\) is smooth of relative dimension
  \(d(r-d)\). Separatedness is checked from the diagonal
  \(\Delta_{I,J} \subset U^I \times U^J\) cut out by \(X^J_I X^I - X^J = 0\).
  Properness follows from the valuative criterion for DVRs: given
  \(\varphi : \Spec K \to \mathrm{Gr}(r,d)\) over a DVR \(R\) with fraction
  field \(K\), factoring through some chart \(U^I\) via a ring map
  \(f : \mathbb{Z}[X^I] \to K\), choose \(J\) minimising the valuation
  \(\nu(f(P^I_J))\) of the minor; on the chart \(U^J\) all entries have
  non-negative valuation, so \(f\) extends uniquely to \(\Spec R\).

  The universal quotient \(u : \mathcal{O}_{\mathrm{Gr}(r,d)}^{\oplus r}
  \twoheadrightarrow \mathcal{U}\) is given on \(U^I\) by the matrix \(X^I\)
  and glued by \(g_{I,J} = (X^I_J)^{-1} \in GL_d(U^I \cap U^J)\). The
  determinant line bundle \(\det\mathcal{U}\) has transition functions
  \(\det g_{I,J} = 1/P^I_J\); the global sections
  \(\sigma_I \in \Gamma(\mathrm{Gr}(r,d), \det\mathcal{U})\) defined by
  \(\sigma_I|_{U^J} = P^J_I\) form a base-point-free linear system separating
  points over \(\Spec\mathbb{Z}\), giving an embedding into
  \(\mathbb{P}^{\binom{r}{d}-1}_{\mathbb{Z}}\), which is closed by
  properness. Representability of \(\mathrm{Grass}(r,d)\) is then verified
  on the universal pair \((\mathrm{Gr}(r,d), u)\): a quotient
  \(q : \mathcal{O}_T^{\oplus r} \twoheadrightarrow \mathcal{F}\) of
  rank \(d\) on \(T\) is everywhere locally trivialised; on the locus where
  the columns indexed by \(I\) are a basis the matrix of \(q\) has the form
  \(X^I\), defining the classifying morphism \(T \to U^I \subset
  \mathrm{Gr}(r,d)\). Base change from \(\mathbb{Z}\) to a general noetherian
  base \(S\) with a vector bundle \(V\) is direct: trivialise \(V\) on an open
  cover and glue. For non-trivial \(V\) the same construction produces
  \(\mathrm{Gr}_S(V, d)\) with Pl\"ucker embedding into
  \(\mathbb{P}_S(\bigwedge^d V)\).
\end{proof}

\section{Representability of the Quot scheme}
\label{sec:quot_representable}

\begin{theorem}\leanok
[Representability of the Quot functor]
  \label{thm:quot_representable}
  \lean{AlgebraicGeometry.Scheme.QuotScheme}
  \uses{def:quot_functor, lem:quot_boundedness, lem:quot_alpha_injective,
    lem:quot_reduction_to_pi_star_W, lem:quot_valuative_criterion,
    thm:grassmannian_representable, thm:flattening_stratification_exists,
    thm:flat_base_change_cohomology}

  % SOURCE: [Nitsure], \S 5, "Main Existence Theorems"
  % (read from references/nitsure-hilbert-quot-src/nitsure-hilbert-quot.tex,
  % L2220-L2273); cf.\ Grothendieck, FGA TDTE-IV.
  % SOURCE QUOTE:
  % "Let \(S\) be a noetherian scheme, \(\pi: X\to S\) a projective morphism,
  %  and \(L\) a relatively very ample line bundle on \(X\). Then for any
  %  coherent \({\mathcal O}_X\)-module \(E\) and any polynomial \(\Phi \in \Q[\lambda]\),
  %  the functor \(\Quot_{E/X/S}^{\Phi,L}\) is representable
  %  by a projective \(S\)-scheme \(\quot_{E/X/S}^{\Phi,L}\)."
  \textit{Source: [Nitsure], \S 5, Theorem (Grothendieck), Theorem
  (Altman--Kleiman) (FGA Explained Ch.~5); cf.\ Grothendieck, FGA TDTE-IV.}
  Let \(S\) be a noetherian scheme, \(\pi : X \to S\) a projective morphism,
  \(L\) a relatively very ample line bundle on \(X\), \(E\) a coherent
  \(\mathcal{O}_X\)-module, and \(\Phi \in \mathbb{Q}[\lambda]\). Then the
  functor \(\Quot^{\Phi,L}_{E/X/S}\) of \cref{def:quot_functor} is
  representable by a projective \(S\)-scheme \(\Quot^{\Phi,L}_{E/X/S}\),
  equipped with a universal quotient
  \[
    q^{\mathrm{univ}} :
    \pi^*_{\Quot} E \twoheadrightarrow \mathcal{F}^{\mathrm{univ}}
    \quad \text{on } X \times_S \Quot^{\Phi,L}_{E/X/S},
  \]
  flat over \(\Quot^{\Phi,L}_{E/X/S}\) with constant Hilbert polynomial
  \(\Phi\) on every fiber. The Hilbert scheme is the special case
  \(E = \mathcal{O}_X\): \(\Hilb^{\Phi,L}_{X/S} = \Quot^{\Phi,L}_{\mathcal{O}_X/X/S}\)
  represents the functor of \(T\)-flat closed subschemes
  \(Y \subset X_T\) with fiberwise Hilbert polynomial \(\Phi\).
\end{theorem}

% SOURCE QUOTE PROOF: this proof is long (~5 named sub-sections in [Nitsure]
% \S 5). Each sub-step is recorded with its own SOURCE QUOTE in the
% sub-lemmas below. The structural skeleton of the proof is restated in
% project notation in the \begin{proof} block.
\begin{proof}
  
  The proof has four steps; each step is a sub-lemma in the next
  sub-sections.
  \begin{enumerate}
    \item \textbf{(Boundedness, \cref{lem:quot_boundedness}.)} Pick the
      Castelnuovo--Mumford regularity bound \(m = m(\mathrm{rank}\,V,
      \mathrm{rank}\,W, \Phi)\) uniformly across all fibers (by Mumford's
      \(m\)-regularity theorem applied to \(E_s\) and to all coherent
      quotients of \(E_s\) with Hilbert polynomial \(\Phi\)). For \(r \ge m\)
      and any \(T \to S\) and any \(T\)-flat quotient
      \(q : E_T \twoheadrightarrow \mathcal{F}\) with Hilbert polynomial
      \(\Phi\), the higher direct images \(R^i (\pi_T)_* \mathcal{F}(r)\)
      and \(R^i (\pi_T)_* E_T(r)\) and \(R^i (\pi_T)_* \mathcal{G}(r)\)
      vanish for \(i \ge 1\) (where \(\mathcal{G} = \ker q\)), and the
      direct images \((\pi_T)_* \mathcal{F}(r)\), \((\pi_T)_* E_T(r)\),
      \((\pi_T)_* \mathcal{G}(r)\) are locally free of ranks fixed by the
      data; the canonical maps \((\pi_T)^* (\pi_T)_*(-)(r) \to (-)(r)\)
      are surjective. Cohomology-and-base-change (\cref{thm:flat_base_change_cohomology},
      Stacks tag 02KH) propagates these properties across \(T\)-base change.
    \item \textbf{(Grassmannian embedding, \cref{lem:quot_alpha_injective}.)}
      Reduce by \cref{lem:quot_reduction_to_pi_star_W} to the case
      \(X = \mathbb{P}_S(V)\), \(E = \pi^* W\) with \(V, W\) locally free
      on \(S\), \(L = \mathcal{O}_{\mathbb{P}(V)}(1)\). Fix \(r \ge m\).
      The Grassmannian embedding
      \[
        \alpha : \Quot^{\Phi,L}_{E/X/S} \longrightarrow
        \mathrm{Grass}\!\left(
          W \otimes_{\mathcal{O}_S} \mathrm{Sym}^r V,\; \Phi(r)
        \right)
      \]
      sends \(\langle \mathcal{F}, q\rangle\) on \(T\) to the locally free
      quotient \((\pi_T)_*(q(r)) : (\pi_T)_* E_T(r) \to (\pi_T)_* \mathcal{F}(r)\).
      This is injective on \(T\)-valued points: from the map \(T \to
      \mathrm{Gr}_S(W \otimes \mathrm{Sym}^r V, \Phi(r))\) one recovers
      the kernel-pull-back \((\pi_T)^* (\pi_T)_* \mathcal{G}(r) \to E_T(r)\),
      whose cokernel is \(q(r) : E_T(r) \to \mathcal{F}(r)\);
      twisting back by \(-r\) recovers \(q\).
    \item \textbf{(Locally closed in the Grassmannian via flattening
      stratification.)} Let \(T_0 = \mathrm{Gr}_S(W \otimes \mathrm{Sym}^r V,
      \Phi(r))\) with universal quotient
      \(f : (W \otimes \mathrm{Sym}^r V)_{T_0} \twoheadrightarrow \mathcal{J}\)
      and kernel \(\mathcal{K}\). Form the composite
      \(h : \pi_{T_0}^* \mathcal{K} \hookrightarrow
      \pi_{T_0}^* (W \otimes \mathrm{Sym}^r V) = \pi_{T_0}^* (\pi_{T_0})_* E_{T_0}
      \twoheadrightarrow E_{T_0}\)
      and let \(q := \mathrm{coker}(h) : E_{T_0} \to \mathcal{F}\). By the
      flattening stratification (\cref{thm:flattening_stratification_exists}
      of \cref{chap:Picard_FlatteningStratification}) applied to \(\mathcal{F}\)
      on \(X_{T_0}\), there exists a locally closed subscheme
      \(T_0^\Phi \subseteq T_0\) such that for any \(\phi : Y \to T_0\), the
      pulled-back cokernel \(\mathcal{F}_Y\) is \(Y\)-flat with fiberwise
      Hilbert polynomial \(\Phi\) iff \(\phi\) factors through \(T_0^\Phi\).
      The pair \((T_0^\Phi, \text{universal cokernel})\) represents
      \(\Quot^{\Phi,L}_{E/X/S}\).
    \item \textbf{(Closed, hence projective,
      \cref{lem:quot_valuative_criterion}.)} The functor
      \(\Quot^{\Phi,L}_{E/X/S}\) satisfies the valuative criterion of
      properness with respect to DVRs (Exercise (7) of \S 1, proved
      using that flatness over a DVR is torsion-freeness). Combined
      with the locally closed embedding of step 3 into the projective
      \(\mathrm{Gr}_S \hookrightarrow \mathbb{P}_S(\bigwedge^{\Phi(r)}
      (W \otimes \mathrm{Sym}^r V))\), this upgrades to a \emph{closed}
      embedding. Hence \(\Quot^{\Phi,L}_{E/X/S}\) is a projective \(S\)-scheme.
  \end{enumerate}
  Finally, by \cref{lem:quot_reduction_to_pi_star_W}, the general case
  (arbitrary \(X\) projective over \(S\) and arbitrary coherent \(E\)) follows
  from the case \(X = \mathbb{P}(V)\), \(E = \pi^* W\): pick a closed
  embedding \(X \subset \mathbb{P}(V)\) and write \(E\) as a coherent quotient
  of \(\pi^*(W)(\nu)\) for some vector bundle \(W\) on \(S\) and integer \(\nu\);
  the Quot scheme of the original data is then a closed subscheme of the
  Quot scheme of the universal data, by the closed-embedding clause of
  \cref{lem:quot_reduction_to_pi_star_W}.
\end{proof}

\section{Sub-lemmas of the construction}
\label{sec:quot_sub_lemmas}

\begin{lemma}[Reduction to the universal case]
  \label{lem:quot_reduction_to_pi_star_W}
  \lean{AlgebraicGeometry.TODO.quotReductionToPiStarW}
  \uses{def:quot_functor}

  % SOURCE: [Nitsure], \S 5, sub-section "Reduction to the case of
  % \(\Quot_{\pi^*W/\PP(V)/S}^{\Phi,L}\)"
  % (read from references/nitsure-hilbert-quot-src/nitsure-hilbert-quot.tex,
  % L2294-L2332).
  % SOURCE QUOTE:
  % "It is enough to prove Theorem REF in the special case
  %  that \(X = \PP(V)\) and
  %  \(E= \pi^*(W)\) where \(V\) and \(W\) are vector bundles on \(S\), as
  %  a consequence of the next lemma.
  %
  %  \begin{lemma}
  %  (i) Let \(\nu\) be any integer.
  %  Then tensoring by \(L^{\nu}\) gives an isomorphism of
  %  functors from \(\Quot_{E/X/S}^{\Phi,L}\) to
  %  \(\Quot_{E(\nu)/X/S}^{\Psi,L}\) where the polynomial
  %  \(\Psi \in \Q[\lambda]\) is defined by \(\Psi(\lambda) = \Phi(\lambda + \nu)\).
  %
  %  (ii) Let \(\phi : E \to G\) be a surjective homomorphism of
  %  coherent sheaves on \(X\). Then the corresponding natural transformation
  %  \(\Quot_{G/X/S}^{\Phi,L} \to \Quot_{E/X/S}^{\Phi,L}\)
  %  is a closed embedding.
  %  \end{lemma}"
  \textit{Source: [Nitsure], \S 5 (FGA Explained Ch.~5).}
  Let \(S\), \(X\), \(L\), \(E\), \(\Phi\) be as in \cref{def:quot_functor}.
  \begin{enumerate}
    \item (Twist-shift.) For \(\nu \in \mathbb{Z}\), tensoring by
      \(L^{\otimes \nu}\) gives an isomorphism of functors
      \(\Quot^{\Phi,L}_{E/X/S} \xrightarrow{\sim}
      \Quot^{\Psi,L}_{E(\nu)/X/S}\) with \(\Psi(\lambda) = \Phi(\lambda + \nu)\).
    \item (Surjection induces closed embedding.) Given a surjection
      \(\phi : E \twoheadrightarrow G\) of coherent sheaves on \(X\), the
      induced natural transformation
      \(\Quot^{\Phi,L}_{G/X/S} \to \Quot^{\Phi,L}_{E/X/S}\) is a closed
      embedding of functors.
  \end{enumerate}
  Consequently, if \(\Quot^{\Phi,L}_{\pi^* W / \mathbb{P}(V)/S}\) is
  representable for all \(V\), \(W\) vector bundles on \(S\), then so is
  \(\Quot^{\Phi,L}_{E/X/S}\) for every \(X \subset \mathbb{P}(V)\) closed and
  every coherent \(E\) on \(X\) that is a quotient of \(\pi^*(W)(\nu)|_X\) for
  some \(W\) and \(\nu\).
\end{lemma}

% SOURCE QUOTE PROOF: "The statement (i) is obvious. The statement (ii) just
% says that given any locally noetherian scheme \(T\) and a family
% \(\langle \F,q\rangle \in \Quot_{E/X/S}^{\Phi,L}(T)\), there exists
% a closed subscheme \(T' \subset T\) with the following universal property:
% for any locally noetherian scheme \(U\) and a morphism \(f: U\to T\),
% the pulled back homomorphism of \({\mathcal O}_{X_U}\)-modules \(q_U : E_U \to \F_U\)
% factors via the pulled back homomorphism \(\phi_U : E_U \to G_U\)
% if and only if \(U \to T\) factors via \(T'\hra T\).
% This is satisfied by taking \(T'\) to be the vanishing scheme
% for the composite homomorphism \(\ker(\phi) \hra  E \stackrel{q}{\to} \F\)
% of coherent sheaves on \(X_T\) (see Remark REF),
% which makes sense here as both \(\ker(\phi)\) and \(\F\) are coherent
% on \(X_T\) and \(\F\) is flat over \(T\)."
\begin{proof}
  
  (i) is direct from the definitions: tensoring with \(L^{\otimes \nu}\) is
  an autoequivalence of coherent sheaves on \(X\) (and on every \(X_T\)),
  preserving surjectivity, \(T\)-flatness, and properness of support, and
  shifting the Hilbert polynomial by \(\nu\) in the variable.

  (ii) Given a family \(\langle \mathcal{F}, q\rangle\) on \(T\) in
  \(\Quot^{\Phi,L}_{E/X/S}\), the locus where \(q\) factors through
  \(\phi : E \twoheadrightarrow G\) is cut out by the vanishing of the
  composite \(\ker(\phi)|_{X_T} \hookrightarrow E_T \xrightarrow{q}
  \mathcal{F}\). Because \(\mathcal{F}\) is \(T\)-flat, the schematic image of
  this composite descends to a closed subscheme \(T' \subset T\) such that
  \(f : U \to T\) factors through \(T'\) iff \(q_U\) factors through \(\phi_U\).
  Hence \(\Quot^{\Phi,L}_{G/X/S} \hookrightarrow \Quot^{\Phi,L}_{E/X/S}\)
  is a closed embedding of subfunctors.
\end{proof}

\begin{lemma}[Uniform Castelnuovo--Mumford regularity bound]
  \label{lem:quot_boundedness}
  \lean{AlgebraicGeometry.TODO.quotBoundedness}
  \uses{thm:flat_base_change_cohomology}

  % SOURCE: [Nitsure], \S 5, sub-section "Use of \(m\)-Regularity"
  % (read from references/nitsure-hilbert-quot-src/nitsure-hilbert-quot.tex,
  % L2346-L2404). Builds on [Nitsure] \S 2 (Castelnuovo--Mumford regularity)
  % and \S 3 (semi-continuity and base-change), cited internally as
  % Theorem~"Mumford" and Theorem~"Base-change for flat sheaves".
  % SOURCE QUOTE:
  % "We consider the sheaf \(E = \pi^*(W)\) on \(X=\PP(V)\) where
  %  \(V\) is a vector bundle on \(S\),
  %  and take \(L = {\mathcal O}_{\PP(V)}(1)\).
  %  For any field \(k\) and a \(k\)-valued point
  %  \(s\) of \(S\), we have an isomorphism
  %  \(\PP(V)_s \simeq \PP^n_k\) where \(n=\rank(V)-1\),
  %  and the restricted sheaf \(E_s\) on \(\PP(V)_s\) is isomorphic to
  %  \(\oplus^p{\mathcal O}_{\PP(V)_s}\) where \(p=\rank(W)\).
  %  It follows from Theorem REF
  %  that given any \(\Phi \in \Q[\lambda]\), there exists an integer
  %  \(m\) which depends only on \(\rank(V)\), \(\rank(W)\) and \(\Phi\),
  %  such that for any field \(k\) and a \(k\)-valued point \(s\) of \(S\),
  %  the sheaf \(E_s\) on \(\PP(V)_s\) is \(m\)-regular, and for any
  %  coherent quotient \(q: E_s \to \F\) on \(\PP(V)_s\) with
  %  Hilbert polynomial \(\Phi\), the sheaf \(\F\) and the kernel sheaf
  %  \(\G\subset E_s\) of \(q\) are both \(m\)-regular."
  \textit{Source: [Nitsure], \S 5, "Use of \(m\)-Regularity" (FGA Explained
  Ch.~5).}
  Let \(X = \mathbb{P}_S(V)\), \(E = \pi^* W\) with \(V\), \(W\) vector bundles on
  \(S\) of ranks \(n+1\), \(p\), \(L = \mathcal{O}_{\mathbb{P}(V)}(1)\), and
  \(\Phi \in \mathbb{Q}[\lambda]\). There exists an integer
  \(m = m(n, p, \Phi)\), depending only on \(n\), \(p\), and \(\Phi\), such that
  for every geometric point \(s\) of \(S\) and every coherent quotient
  \(q : E_s \twoheadrightarrow \mathcal{F}\) with Hilbert polynomial
  \(\Phi\), the sheaves \(E_s\), \(\mathcal{F}\), and \(\ker q\) on \(X_s\) are all
  \(m\)-regular in the sense of Castelnuovo--Mumford. Consequently, for
  every \(T \to S\) and every \(T\)-flat quotient
  \(q : E_T \twoheadrightarrow \mathcal{F}\) on \(X_T\) with Hilbert
  polynomial \(\Phi\) and every \(r \ge m\), the conclusions \textbf{(*)} and
  \textbf{(**)} of \cite[\S 5, "Use of \(m\)-Regularity"]{nitsure-hilbert-quot}
  hold: the direct images \((\pi_T)_* \mathcal{F}(r)\), \((\pi_T)_* E_T(r)\),
  \((\pi_T)_* \mathcal{G}(r)\) are locally free of fixed ranks; the
  canonical surjections \((\pi_T)^* (\pi_T)_*(-)(r) \twoheadrightarrow (-)(r)\)
  hold; and the higher direct images \(R^i (\pi_T)_* (-)(r)\) vanish for
  \(i \ge 1\).
\end{lemma}

% SOURCE QUOTE PROOF: see [Nitsure], \S 5, "Use of \(m\)-Regularity"
% (references/nitsure-hilbert-quot-src/nitsure-hilbert-quot.tex, L2346-L2404).
% The proof invokes the m-regularity theorem of [Nitsure] \S 2
% (Theorem "Mumford") to extract the uniform bound m, and the
% base-change-for-flat-sheaves theorem of [Nitsure] \S 3 (cited as
% "Base-change for flat sheaves") to propagate the conclusions across \(T\).
\begin{proof}
  
  Each geometric fiber \(E_s\) is a trivial bundle on \(\mathbb{P}^n_{\kappa(s)}\)
  of rank \(p\), so \(E_s\) is \(0\)-regular. The Castelnuovo--Mumford
  \(m\)-regularity theorem of [Nitsure] \S 2 produces, for any fixed Hilbert
  polynomial \(\Phi\), a uniform integer \(m = m(n, p, \Phi)\) above which
  every coherent quotient of \(E_s\) with Hilbert polynomial \(\Phi\) -- and
  its kernel -- is \(m\)-regular. The Castelnuovo Lemma of [Nitsure] \S 2
  converts \(m\)-regularity into vanishing of \(H^i((-) (r))\) for \(i \ge 1\)
  and \(r \ge m\), and into global generation of \(H^0((-)(r))\). To
  propagate the fiber-wise statement to families over \(T\), apply the
  flat-base-change theorem (\cref{thm:flat_base_change_cohomology}, Stacks
  tag 02KH) along \(T \to S\): cohomology of flat families commutes with
  base change, so the direct images \((\pi_T)_*(-)(r)\) are locally free of
  the rank dictated by the fiber Hilbert polynomials, and the canonical
  maps \((\pi_T)^* (\pi_T)_*(-)(r) \to (-)(r)\) inherit fiberwise
  surjectivity.
\end{proof}

\begin{lemma}[Grassmannian embedding \(\alpha\) is injective]
  \label{lem:quot_alpha_injective}
  \lean{AlgebraicGeometry.TODO.quotAlphaInjective}
  \uses{lem:quot_boundedness}

  % SOURCE: [Nitsure], \S 5, sub-section "Embedding Quot into Grassmannian"
  % (read from references/nitsure-hilbert-quot-src/nitsure-hilbert-quot.tex,
  % L2409-L2451).
  % SOURCE QUOTE:
  % "We now fix a positive integer \(r\) such that \(r\ge m\).
  %  Note that the rank of \({\pi_T}_*\F(r)\) is \(\Phi(r)\)
  %  and \(\pi_*E(r) = W \otimes_{{\mathcal O}_S}\sym^r V\).
  %  Therefore the surjective homomorphism
  %  \({\pi_T}_*E_T(r) \to {\pi_T}_*\F(r)\)
  %  defines an element of the set
  %  \(\Grass( W \otimes_{{\mathcal O}_S}\sym^r V , \Phi(r))(T)\).
  %  We thus get a morphism of functors
  %  \(\)\alpha : \Quot^{\Phi,L}_{E/X/S} \to
  %  \Grass( W \otimes_{{\mathcal O}_S}\sym^r V, \Phi(r))\(\)
  %  It associates to \(q: E_T \to \F\) the quotient
  %  \({\pi_T}_*(q(r)) : {\pi_T}_*E_T(r) \to {\pi_T}_*\F(r)\).
  %
  %  The above morphism \(\alpha\) is injective because the quotient
  %  \(q: E_T \to \F\) can be recovered from
  %  \({\pi_T}_*(q(r)) : {\pi_T}_*E_T(r) \to {\pi_T}_*\F(r)\)
  %  as follows."
  \textit{Source: [Nitsure], \S 5, "Embedding Quot into Grassmannian"
  (FGA Explained Ch.~5).}
  In the universal setup \(X = \mathbb{P}_S(V)\), \(E = \pi^* W\), fix
  \(r \ge m\) from \cref{lem:quot_boundedness}. The natural transformation
  \[
    \alpha :
    \Quot^{\Phi,L}_{E/X/S} \longrightarrow
    \mathrm{Grass}\!\left(
      W \otimes_{\mathcal{O}_S} \mathrm{Sym}^r V,\; \Phi(r)
    \right),
    \qquad
    \langle \mathcal{F}, q\rangle \longmapsto
    \langle (\pi_T)_*\mathcal{F}(r),\, (\pi_T)_*(q(r))\rangle
  \]
  is well-defined (the target is locally free of rank \(\Phi(r)\), and the
  canonical isomorphism \(\pi_* E(r) \cong W \otimes_{\mathcal{O}_S}
  \mathrm{Sym}^r V\) gives the source) and \emph{injective} on \(T\)-points:
  the data of the locally free quotient at level \(r\) determines
  \(q : E_T \to \mathcal{F}\).
\end{lemma}

\begin{proof}
  
  Suppose \(\alpha(T)\) sends both \(\langle \mathcal{F}_1, q_1\rangle\) and
  \(\langle \mathcal{F}_2, q_2\rangle\) to the same \(T\)-point of the
  Grassmannian. Then the pulled-back kernel
  \((\pi_T)^* (\pi_T)_* \mathcal{G}_i(r) \to (\pi_T)^* (\pi_T)_* E_T(r)
  = \pi_{T_0}^* (W \otimes \mathrm{Sym}^r V)_T\)
  is the same. The diagram \textbf{(**)} of
  \cref{lem:quot_boundedness} gives a right-exact sequence
  \[
    (\pi_T)^* (\pi_T)_* \mathcal{G}_i(r) \xrightarrow{h}
    E_T(r) \xrightarrow{q_i(r)} \mathcal{F}_i(r) \to 0,
  \]
  so \(q_i(r)\) is recovered as the cokernel of \(h\); twisting by \(-r\)
  recovers \(q_i\), giving \(q_1 = q_2\) and hence
  \(\langle \mathcal{F}_1, q_1\rangle = \langle \mathcal{F}_2, q_2\rangle\).
\end{proof}

\begin{lemma}[Valuative criterion of properness for \(\Quot\)]
  \label{lem:quot_valuative_criterion}
  \lean{AlgebraicGeometry.TODO.quotValuativeCriterion}
  % SOURCE: [Nitsure], \S 5, sub-section "Valuative Criterion for Properness"
  % (read from references/nitsure-hilbert-quot-src/nitsure-hilbert-quot.tex,
  % L2510-L2542). Original reference: EGA IV(2) 2.8.1, cited there.
  % SOURCE QUOTE:
  % "The functor \(\Quot_{E/X/S}^{\Phi,L}\) satisfies the following valuative
  %  criterion for properness over \(S\): given any discrete valuation ring
  %  \(R\) over \(S\) with quotient field \(K\), the restriction map
  %  \(\)\Quot_{E/X/S}^{\Phi,L}(\spec R) \to \Quot_{E/X/S}^{\Phi,L}(\spec K)\(\)
  %  is bijective. This can be seen as follows.
  %  Given any coherent quotient
  %  \(q: E_K \to \F\) on \(X_R\) which defines an element \(\langle \F, q\rangle\)
  %  of \(\Quot_{E/X/S}^{\Phi,L}(\spec K)\).
  %  Let \(\ov{\F}\) be the image of the composite homomorphism
  %  \(E_R \to j_*(E_K) \to j_*\F\) where \(j : X_K \hra X_R\) is the open
  %  inclusion. Let \(\ov{q} : E_R\to \ov{\F}\) be the induced surjection.
  %  Then we leave it to the reader to verify that
  %  \(\langle \F, q\rangle\) is an element of \(\Quot_{E/X/S}^{\Phi,L}(\spec R)\)
  %  which maps to \(\langle \F, q\rangle\), and is the unique such element.
  %  (Use the basic fact that being flat over a dvr
  %  is the same as being torsion-free.)"
  \textit{Source: [Nitsure], \S 5, "Valuative Criterion for Properness"
  (FGA Explained Ch.~5); EGA IV(2) 2.8.1.}
  The functor \(\Quot^{\Phi,L}_{E/X/S}\) satisfies the valuative criterion
  for properness over \(S\): for every discrete valuation ring \(R\) with
  fraction field \(K\) and every morphism \(\Spec R \to S\), the restriction
  map
  \[
    \Quot^{\Phi,L}_{E/X/S}(\Spec R)
    \longrightarrow
    \Quot^{\Phi,L}_{E/X/S}(\Spec K)
  \]
  is bijective.
\end{lemma}

\begin{proof}
  
  Given a \(\Spec K\)-point \(\langle \mathcal{F}_K, q_K\rangle\) on \(X_K\),
  let \(j : X_K \hookrightarrow X_R\) be the open inclusion. Define
  \(\overline{\mathcal{F}}\) as the image of the composite
  \(E_R \to j_* E_K \to j_* \mathcal{F}_K\), with induced surjection
  \(\overline{q} : E_R \to \overline{\mathcal{F}}\). The sheaf
  \(\overline{\mathcal{F}}\) is torsion-free over \(R\) (being a subsheaf of
  \(j_* \mathcal{F}_K\)), hence \(R\)-flat (flatness over a DVR
  \(=\) torsion-freeness), and restricts on the generic fiber to
  \(\mathcal{F}_K\). The pair \((\overline{\mathcal{F}}, \overline{q})\) is
  the unique extension of \(\langle \mathcal{F}_K, q_K\rangle\) to a
  \(\Spec R\)-point (uniqueness: any other extension agrees on the dense
  open \(X_K \subset X_R\) and is torsion-free, hence equals the image
  construction). The Hilbert polynomial is constant on connected
  \(\Spec R\) by flatness.
\end{proof}

\section{Cohomology and base change}
\label{sec:cohomology_base_change}

The Quot construction uses cohomology-and-base-change in two places: the
boundedness step (\cref{lem:quot_boundedness}) and the Grassmannian
embedding (\cref{lem:quot_alpha_injective}). We record the relevant
statement as a named theorem so the Lean encoding can cite it directly.

\begin{theorem}\leanok
[Flat base change of cohomology]
  \label{thm:flat_base_change_cohomology}
  \lean{AlgebraicGeometry.flatBaseChangeCohomology}
  \uses{thm:flat_base_change_higher, thm:flat_base_change_pushforward,
    lem:higher_direct_image_quasi_coherent}
  % SOURCE: [Stacks Project], tag 02KH (lemma-flat-base-change-cohomology)
  % (read from references/stacks-coherent.tex, L947-L970).
  % SOURCE QUOTE:
  % "Consider a cartesian diagram of schemes
  %  \(\)
  %  \xymatrix{
  %  X' \ar[d]_{f'} \ar[r]_{g'} & X \ar[d]^f \\
  %  S' \ar[r]^g & S
  %  }
  %  \(\)
  %  Let \(\mathcal{F}\) be a quasi-coherent \(\mathcal{O}_X\)-module
  %  with pullback \(\mathcal{F}' = (g')^*\mathcal{F}\).
  %  Assume that \(g\) is flat and that \(f\) is quasi-compact and quasi-separated.
  %  For any \(i \geq 0\)
  %  \begin{enumerate}
  %  \item the base change map of
  %  Cohomology, Lemma REF
  %  is an isomorphism
  %  \(\)
  %  g^*R^if_*\mathcal{F} \longrightarrow R^if'_*\mathcal{F}',
  %  \(\)
  %  \item if \(S = \Spec(A)\) and \(S' = \Spec(B)\), then
  %  \(H^i(X, \mathcal{F}) \otimes_A B = H^i(X', \mathcal{F}')\).
  %  \end{enumerate}"
  \textit{Source: [Stacks Project], tag 02KH (lemma-flat-base-change-cohomology).}
  Let
  \[
    \begin{array}{ccc}
      X' & \xrightarrow{g'} & X \\
      \downarrow f' &  & \downarrow f \\
      S' & \xrightarrow{g} & S
    \end{array}
  \]
  be a cartesian diagram of schemes with \(g\) flat and \(f\) quasi-compact and
  quasi-separated. Let \(\mathcal{F}\) be a quasi-coherent
  \(\mathcal{O}_X\)-module with pull-back
  \(\mathcal{F}' = (g')^* \mathcal{F}\). For every \(i \ge 0\):
  \begin{enumerate}
    \item the base-change map
      \(g^* R^i f_* \mathcal{F} \to R^i f'_* \mathcal{F}'\) is an isomorphism;
    \item if \(S = \Spec A\) and \(S' = \Spec B\), then \(H^i(X, \mathcal{F})
      \otimes_A B = H^i(X', \mathcal{F}')\).
  \end{enumerate}
\end{theorem}

\subsection{Project-side base-change substrate}
\label{subsec:quot_project_basechange}

The Lean encoding of \cref{thm:flat_base_change_cohomology} (and the
Quot construction's iso between pushforward and tensor-pullback for the
universal section pattern) relies on four project-side declarations
introduced in \texttt{AlgebraicJacobian/Picard/QuotScheme.lean}. These
are typed-sorry today: each carries a substantive Lean type that
encodes the intended mathematical statement, with the body deferred to
later iters (see \cref{sec:quot_lean_encoding} for the
phase-by-phase plan).

\begin{definition}\leanok
[Canonical base-change map on sheaves of modules]
  \label{def:quot_canonical_basechange_map}
  \lean{AlgebraicGeometry.canonicalBaseChangeMap}
  \uses{thm:flat_base_change_cohomology}
  Given a cartesian square of schemes
  \(\square = (g' : X' \to X,\ f' : X' \to S',\ f : X \to S,\ g : S' \to S)\)
  with \(g \circ f' = f \circ g'\) and \(g\) flat, the
  \emph{canonical base-change map} on sheaves of modules is the natural
  transformation \(g'^* f^* \mathcal{F} \longrightarrow f'^* g^* \mathcal{F}\)
  arising from the universal property of pullback. Encoded as a morphism
  of \(\mathrm{ModuleCat}\)-valued sheaf functors on the pulled-back open
  cover; this is the morphism whose iso-status discharges
  \cref{thm:flat_base_change_cohomology}(1) in the present
  finite-locally-free / quasi-coherent-of-finite-type setting.
\end{definition}

\begin{theorem}\leanok
[Stalkwise base-change isomorphism on affines]
  \label{thm:quot_canonical_basechange_app_app_isIso}
  \lean{AlgebraicGeometry.canonicalBaseChangeMap_app_app_isIso}
  \uses{def:quot_canonical_basechange_map, thm:flat_base_change_cohomology, thm:quot_canonical_basechange_app_app_isIso_of_affineCover, thm:quot_canonical_basechange_app_app_isIso_of_isAffineOpen}
  For every affine open \(V \subseteq S'\) and every affine open
  \(U \subseteq X\) lying over \(V\), the local component
  \((\text{canonicalBaseChangeMap}\ \square).\text{app}\ \mathcal{F}.\text{app}\ U\)
  of \cref{def:quot_canonical_basechange_map} is an isomorphism. Proof
  reduces (Stacks 02KH(2)) to the affine-base case: \(H^0\) commutes with
  flat extension of scalars on a quasi-compact quasi-separated affine
  morphism.
\end{theorem}

\begin{theorem}

\leanok
[Affine-base case of the stalkwise base-change isomorphism]
  \label{thm:quot_canonical_basechange_app_app_isIso_of_isAffineBase}
  \lean{AlgebraicGeometry.canonicalBaseChangeMap_app_app_isIso_of_isAffineOpen_of_isAffineBase}
  \uses{def:quot_pullback_app_isoTensor, lem:pushforward_isQuasicoherent, lem:pullback_tildeIso}
  Assume in addition that the base \(S\) is affine, that \(f\) is
  quasi-compact and quasi-separated, that \(g\) is flat, and that
  \(\mathcal{F}\) is quasi-coherent. Then for every affine open
  \(U \subseteq S'\) the section of the canonical base-change map
  (\cref{def:quot_canonical_basechange_map}) over \(U\) is an isomorphism.
  Taking \(V = S\) (affine) as the compatible base open, the statement
  reduces to the algebraic flat base-change identity: the global sections
  \(\Gamma(S', U)\) are flat over \(\Gamma(S, S)\), and the
  tensor-presentation of the pulled-back sections
  (\cref{def:quot_pullback_app_isoTensor}, applied to the pushforward
  \((\pi_f)_* \mathcal{F}\), which stays quasi-coherent by
  \cref{lem:pushforward_isQuasicoherent}) matches the base-change map under
  the Stacks~01HQ identification of \cref{lem:pullback_tildeIso}.
\end{theorem}
\begin{proof} Proved directly in Lean. \end{proof}

\begin{theorem}

\leanok
[Affine-open case of the stalkwise base-change isomorphism]
  \label{thm:quot_canonical_basechange_app_app_isIso_of_isAffineOpen}
  \lean{AlgebraicGeometry.canonicalBaseChangeMap_app_app_isIso_of_isAffineOpen}
  \uses{thm:quot_canonical_basechange_app_app_isIso_of_isAffineBase, lem:quot_pushforward_pullback_section_eq}
  With \(f\) quasi-compact and quasi-separated, \(g\) flat, and
  \(\mathcal{F}\) quasi-coherent (but \(S\) now arbitrary), the section of
  the canonical base-change map (\cref{def:quot_canonical_basechange_map})
  over an affine open \(U \subseteq S'\) is an isomorphism. The substantive
  algebraic content is the affine-base specialization
  (\cref{thm:quot_canonical_basechange_app_app_isIso_of_isAffineBase}); the
  passage from a general base \(S\) to the affine case is a base-side
  Mayer-Vietoris descent along a finite affine cover of \(S\), using the
  definitional section identity
  \cref{lem:quot_pushforward_pullback_section_eq} on each piece.
\end{theorem}
\begin{proof} Proved directly in Lean. \end{proof}

\begin{theorem}

\leanok
[Open-cover descent for the stalkwise base-change isomorphism]
  \label{thm:quot_canonical_basechange_app_app_isIso_of_affineCover}
  \lean{AlgebraicGeometry.canonicalBaseChangeMap_app_app_isIso_of_affineCover}
  \uses{def:quot_canonical_basechange_map}
  Suppose the section of the canonical base-change map
  (\cref{def:quot_canonical_basechange_map}) is an isomorphism over
  \emph{every} affine open of \(S'\). Then it is an isomorphism over every
  open \(U \subseteq S'\). This is the standard Mayer-Vietoris descent for a
  morphism of quasi-coherent sheaves on the base: cover \(U\) by affine
  opens, on each of which the section is an isomorphism by hypothesis, and
  glue along the intersections (quasi-compact because \(f\) is
  quasi-separated). It is the descent half of the stalkwise statement
  \cref{thm:quot_canonical_basechange_app_app_isIso}, whose Lean proof feeds
  the affine-open case
  \cref{thm:quot_canonical_basechange_app_app_isIso_of_isAffineOpen} as the
  affine-open hypothesis.
\end{theorem}
\begin{proof} Proved directly in Lean. \end{proof}

\begin{lemma}

\leanok
[Definitional section identity for pushforward of a pullback]
  \label{lem:quot_pushforward_pullback_section_eq}
  \lean{AlgebraicGeometry.pushforward_pullback_section_eq_pullback_section}
  \uses{def:quot_canonical_basechange_map}
  Let \(f' : X' \to S'\) and \(g' : X' \to X\) be morphisms of schemes,
  \(\mathcal{F}\) a sheaf of \(\mathcal{O}_X\)-modules, and \(U \subseteq S'\)
  an open. The sections over \(U\) of the pushforward along \(f'\) of the
  pullback along \(g'\) of \(\mathcal{F}\) coincide \emph{definitionally}
  with the sections of the pullback \(g'^* \mathcal{F}\) over the preimage
  \(f'^{-1}(U)\):
  \[
    \Gamma\bigl(U,\ f'_* (g'^* \mathcal{F})\bigr)
    \;=\;
    \Gamma\bigl(f'^{-1}(U),\ g'^* \mathcal{F}\bigr).
  \]
  This records step~(3) of the affine-open construction
  (\cref{thm:quot_canonical_basechange_app_app_isIso_of_isAffineOpen}): the
  target of the canonical base-change map
  (\cref{def:quot_canonical_basechange_map}) is read through this identity.
\end{lemma}
\begin{proof} Proved directly in Lean. \end{proof}

\begin{theorem}\leanok
[Canonical base-change map is a sheaf isomorphism]
  \label{thm:quot_canonical_basechange_isIso}
  \lean{AlgebraicGeometry.canonicalBaseChangeMap_isIso}
  \uses{thm:quot_canonical_basechange_app_app_isIso}
  Globalising the stalk-wise statement
  \cref{thm:quot_canonical_basechange_app_app_isIso}: the natural
  transformation \(g'^* f^* \mathcal{F} \to f'^* g^* \mathcal{F}\) of
  \cref{def:quot_canonical_basechange_map} is an isomorphism of
  quasi-coherent sheaves of modules. This is the sheaf-level form of
  \cref{thm:flat_base_change_cohomology}(1).
\end{theorem}

\begin{definition}\leanok
[Tensor-presentation of the pullback-of-sections functor]
  \label{def:quot_pullback_app_isoTensor}
  \lean{AlgebraicGeometry.Scheme.Modules.pullback_app_isoTensor}
  \uses{def:quot_canonical_basechange_map, thm:quot_pullback_app_isoTensor_isBaseChange}
  Given a morphism \(g : X \to Y\) of schemes, affine opens
  \(U \subseteq Y\), \(V \subseteq X\) with \(g(V) \subseteq U\), and a
  quasi-coherent sheaf of modules \(\mathcal{N}\) on \(Y\), the
  \emph{tensor-presentation iso} packages the canonical iso
  \(\Gamma(V,\, g^* \mathcal{N}) \;\cong\;
   \Gamma(V,\, \mathcal{O}_X) \otimes_{\Gamma(U,\, \mathcal{O}_Y)}
     \Gamma(U,\, \mathcal{N})\)
  as a `LinearEquiv`. This is the project-side rfl-bridge used by the
  consumer \(\text{Quot}\)-functor iso assembly: it discharges the
  pullback-app step of the Tilde-isoTop route (Stacks tag 01HQ / 01I8)
  to a base-change-of-scalars identity. The body factors as
  (Step 1, axiom-clean): the underlying linear map
  \texttt{pullback\_app\_isoTensor\_unitAtV};
  (Step 2, project-side typed sorry): the
  \texttt{IsBaseChange}-packaging step
  \texttt{pullback\_app\_isoTensor\_isBaseChange} reducing to
  Stacks 02KE bijectivity of the canonical map.
\end{definition}

\begin{definition}

\leanok
[Adjunction-unit base linear map at a section]
  \label{def:quot_pullback_app_isoTensor_unitAtV}
  \lean{AlgebraicGeometry.pullback_app_isoTensor_unitAtV}
  Given a morphism \(g : Y \to X\) of schemes, a sheaf of
  \(\mathcal{O}_X\)-modules \(\mathcal{N}\), and an open \(V \subseteq X\),
  the unit of the pullback--pushforward adjunction yields a morphism of
  \(\mathcal{O}_X\)-modules \(\mathcal{N} \to g_*(g^* \mathcal{N})\).
  Evaluating its \(V\)-component gives the \(\Gamma(X, V)\)-linear
  \emph{unit-at-section} map
  \[
    \Gamma(\mathcal{N}, V)
    \;\longrightarrow\;
    \Gamma\bigl(g_*(g^* \mathcal{N}),\ V\bigr).
  \]
  This is the axiom-clean building block (Step~1 of the Tilde-isoTop route)
  underlying the tensor-presentation iso of
  \cref{def:quot_pullback_app_isoTensor}.
\end{definition}
\begin{proof} Proved directly in Lean. \end{proof}

\begin{definition}

\leanok
[Section-level base map for the pullback]
  \label{def:quot_pullback_app_isoTensor_baseMap}
  \lean{AlgebraicGeometry.pullback_app_isoTensor_baseMap}
  \uses{def:quot_pullback_app_isoTensor_unitAtV}
  For a morphism \(g : Y \to X\), a sheaf \(\mathcal{N}\) on \(X\), and
  compatible affine opens \(U \subseteq Y\), \(V \subseteq X\) with
  \(g(U) \subseteq V\), composing the unit-at-section map
  (\cref{def:quot_pullback_app_isoTensor_unitAtV}) with the presheaf
  restriction from \(g^{-1}(V)\) down to \(U\) produces the
  \(\Gamma(X, V)\)-linear \emph{base map}
  \[
    \Gamma(\mathcal{N}, V)
    \;\longrightarrow\;
    \Gamma\bigl(g^* \mathcal{N},\ U\bigr),
  \]
  where the \(\Gamma(X, V)\)-action on the target is restriction of scalars
  along the ring map \(\Gamma(X, V) \to \Gamma(Y, U)\) induced by \(g\). This
  is the candidate linear map whose base-change property
  (\cref{thm:quot_pullback_app_isoTensor_baseMap_isBaseChange}) yields the
  tensor presentation.
\end{definition}
\begin{proof} Proved directly in Lean. \end{proof}

\begin{theorem}

\leanok
[The base map is a base change]
  \label{thm:quot_pullback_app_isoTensor_baseMap_isBaseChange}
  \lean{AlgebraicGeometry.pullback_app_isoTensor_baseMap_isBaseChange}
  \uses{def:quot_pullback_app_isoTensor_baseMap, def:pullback_app_isoTensor_sigma}
  For compatible affine opens \(U \subseteq Y\), \(V \subseteq X\) of a
  morphism \(g : Y \to X\) and a quasi-coherent sheaf \(\mathcal{N}\) on
  \(X\), the base map of \cref{def:quot_pullback_app_isoTensor_baseMap}
  exhibits \(\Gamma(g^* \mathcal{N}, U)\) as the base change of
  \(\Gamma(\mathcal{N}, V)\) along the ring map
  \(\Gamma(X, V) \to \Gamma(Y, U)\); i.e.\ the induced
  \(\Gamma(Y, U) \otimes_{\Gamma(X, V)} \Gamma(\mathcal{N}, V)
   \to \Gamma(g^* \mathcal{N}, U)\) is an isomorphism. The proof feeds the
  packaged section-level linear equivalence and its intertwining property
  (\cref{def:pullback_app_isoTensor_sigma}) into the universal
  characterization of a base change.
\end{theorem}
\begin{proof} Proved directly in Lean. \end{proof}

\begin{theorem}

\leanok
[Existence of the tensor-presentation equivalence]
  \label{thm:quot_pullback_app_isoTensor_isBaseChange}
  \lean{AlgebraicGeometry.pullback_app_isoTensor_isBaseChange}
  \uses{thm:quot_pullback_app_isoTensor_baseMap_isBaseChange}
  In the same affine setup, there exists a \(\Gamma(Y, U)\)-linear
  isomorphism
  \[
    \Gamma\bigl(g^* \mathcal{N},\ U\bigr)
    \;\cong\;
    \Gamma(Y, U) \otimes_{\Gamma(X, V)} \Gamma(\mathcal{N}, V),
  \]
  obtained from the base-change property
  (\cref{thm:quot_pullback_app_isoTensor_baseMap_isBaseChange}) by passing to
  the induced equivalence. This is the existence statement consumed by the
  tensor-presentation iso \cref{def:quot_pullback_app_isoTensor}.
\end{theorem}
\begin{proof} Proved directly in Lean. \end{proof}

\begin{definition}

\leanok
[\(\Sigma\)-pair packaging: section-level LinearEquiv + intertwining]
  \label{def:pullback_app_isoTensor_sigma}
  \lean{AlgebraicGeometry.pullback_app_isoTensor_baseMap_sectionLinearEquiv}
  % NOTE iter-283: replaced the spurious `def:quot_pullback_app_isoTensor` with
  % `def:quot_pullback_app_isoTensor_baseMap` to break a dependency cycle that
  % crashed `leanblueprint web`. In Lean (QuotScheme.lean L867) this Sigma-pair
  % is built from `tildeIso_of_isQuasicoherent_isAffineOpen`/`pullback_tildeIso`
  % and its statement references `pullback_app_isoTensor_baseMap`, NOT the
  % high-level final iso `pullback_app_isoTensor`. The final iso is built FROM
  % this Sigma-pair (def:quot_pullback_app_isoTensor -> thm:..._isBaseChange ->
  % thm:..._baseMap_isBaseChange -> def:pullback_app_isoTensor_sigma), so the
  % former edge was the back-edge closing the cycle. The corrected dependency
  % (on the lower-level baseMap) is both acyclic and accurate.
  \uses{def:quot_pullback_app_isoTensor_baseMap, def:quot_canonical_basechange_map,
        lem:pullback_tildeIso, lem:tildeIso_of_isQuasicoherent_isAffineOpen,
        lem:pullback_of_openImmersion_iso_restrict, lem:pushforward_isQuasicoherent}
  Iter-188 Lane F-introduced \(\Sigma\)-pair packaging declaration that
  bundles BOTH the section-level LinearEquiv
  \(\Gamma(V,(g^*\mathcal N)) \;\simeq_{\Gamma(Y,U)}\; \Gamma(Y,U) \otimes
    \Gamma(N,V)\) AND its intertwining property with the canonical
  base-change map into a single \(\Sigma\)-struct returned via
  \texttt{Nonempty} for \texttt{IsBaseChange.of\_equiv} consumption.
  Concretely:
  \[
    \exists\, (f : \Gamma(V,(\text{pullback}\, g).\mathrm{obj}\,
                     \mathcal N) \;\simeq_{\Gamma(Y,U)}\;
                     \Gamma(X,V) \otimes_{\Gamma(Y,U)} \Gamma(N,V)),\;
      \forall x,\; f(1\otimes x)\;=\;\text{baseMap}\,g\,\mathcal N\,e\,x.
  \]
  This packaging idiom keeps consumer signatures clean and is the project-
  side bridge into Mathlib's \texttt{IsBaseChange.of\_equiv} consumer.
  Body is gated on \cref{lem:pullback_tildeIso} (Stacks 01HQ) plus a
  6-stage Beck-Chevalley intertwining decomposition recorded in
  \cref{sec:quot_iter195_bc_substrate} and mirroring the in-file
  comment block at
  \texttt{AlgebraicJacobian/Picard/QuotScheme.lean:999-1037}:
  \begin{itemize}
    \item \textbf{Stage 1} (closed iter-195 via the \(\Sigma\)-pair
      iso-characterizing identity \texttt{\_step2\_apply} together with
      \texttt{inv\_hom\_id}): rewrites
      \((\text{step2}.\mathrm{inv}.\mathrm{app}\,\top)\bigl(\text{tilde.toOpen}\,TR\,\top\,(1 \otimes x)\bigr)\)
      as
      \(\text{baseMap}\,(\text{Spec.map}\,\varphi)\,(\widetilde{\Gamma(\mathcal N, V)})\,e_1\,(\text{tilde.toOpen}\,\Gamma(\mathcal N, V)\,\top\,x)\).
    \item \textbf{Stage 2} (uses (N1) baseMap naturality of
      \cref{lem:baseMap_pullbackComp_apply} together with the
      \(\Sigma\)-pair \texttt{\_step1\_apply} identity): transports
      \(\text{step1}.\mathrm{inv}\) across
      \(\text{baseMap}\,(\text{Spec.map}\,\varphi)\), producing
      \(\text{baseMap}\,(\text{Spec.map}\,\varphi)\,((\text{pullback}\,\mathtt{\_hV.fromSpec}).\mathrm{obj}\,\mathcal N)\,e_2\,(\text{baseMap}\,\mathtt{\_hV.fromSpec}\,\mathcal N\,e_3\,x)\).
    \item \textbf{Stage 3} (uses (N2) baseMap composition via
      \texttt{pullbackComp} of
      \cref{lem:baseMap_pullback_comp_apply}): collapses the nested
      pair of \texttt{baseMap}'s of Stage 2 into a single
      \(\text{baseMap}\,(\text{Spec.map}\,\varphi \circ \mathtt{\_hV.fromSpec})\,\mathcal N\,e_4\,x\)
      along the composite morphism.
    \item \textbf{Stage 4} (uses (N3) baseMap transport via
      \texttt{pullbackCongr} of
      \cref{lem:baseMap_pullbackCongr_apply}): transports the
      composite of Stage 3 along the propositional equality
      \(\text{Spec.map}\,\varphi \circ \mathtt{\_hV.fromSpec}
      = \mathtt{\_hU.fromSpec} \circ g\) (which holds by definition of
      \texttt{Scheme.Hom.appLE} on \(g\) with \(g(V) \subseteq U\)),
      producing
      \(\text{baseMap}\,(\mathtt{\_hU.fromSpec} \circ g)\,\mathcal N\,e_5\,x\).
    \item \textbf{Stage 5} (uses (N2) again,
      \cref{lem:baseMap_pullback_comp_apply}): re-decomposes the
      composite of Stage 4 into nested \texttt{baseMap}'s along
      \(\mathtt{\_hU.fromSpec}\) and \(g\), exposing
      \(\text{baseMap}\,g\,\mathcal N\,e\,x\) as the inner section.
    \item \textbf{Stage 6} (uses (N4) \texttt{step3} inversion of
      \cref{lem:baseMap_inv_step3_open_immersion}): inverts the outer
      \(\text{baseMap}\,\mathtt{\_hU.fromSpec}\) via \texttt{step3} —
      the transport iso of
      \cref{lem:pullback_of_openImmersion_iso_restrict} for the open
      immersion \(\mathtt{\_hU.fromSpec}\) — leaving
      \(\text{baseMap}\,g\,\mathcal N\,e\,x\), which matches the RHS
      of the Σ-pair intertwining property.
  \end{itemize}
\end{definition}

\subsection{Iter-187 analogist-licensed named substrate gaps}
\label{sec:quot_analogist_gaps}

The iter-187 \texttt{mathlib-analogist quotscheme-isbasechange-tilde}
verdict (see \texttt{analogies/quotscheme-isbasechange-tilde.md})
identified two distinct Mathlib substrate gaps at the pinned commit
\texttt{b80f227}. Both are pinned here as named project-side typed
sorries; closing either is iter-189+ sub-build work.

\begin{lemma}

\leanok
[Pullback of tilde is tilde of base change (Stacks 01HQ)]
  \label{lem:pullback_tildeIso}
  \lean{AlgebraicGeometry.pullback_tildeIso}
  \uses{def:quot_canonical_basechange_map}
  \textit{Source: Stacks~\href{https://stacks.math.columbia.edu/tag/01HQ}{01HQ}
  / \href{https://stacks.math.columbia.edu/tag/0BJ8}{0BJ8} (pullback of
  quasi-coherent modules along a Spec morphism).}
  For a ring homomorphism \(\varphi : A \to B\) and an \(A\)-module
  \(M\), there is a canonical iso between the pullback of the
  tilde-sheaf and the tilde of the base-changed module:
  \[
    (\mathtt{Scheme.Modules.pullback}\, (\mathrm{Spec.map}\, \varphi)).\mathrm{obj}\,
      (\widetilde{M})
    \;\cong\;
    \widetilde{B \otimes_A M}.
  \]
  Mathlib \texttt{b80f227} ships
  \texttt{PullbackFree.lean:122 pullbackObjFreeIso} for \emph{free} sheaves
  only — the general-module form (Stacks 01HQ) is unowned.
\end{lemma}
\begin{proof}
  \uses{def:quot_canonical_basechange_map}
  Substantive body (Stacks 01HQ): construct via the
  \texttt{tilde.adjunction} naturality square; the construction passes
  through Spec-restriction transport.

  % NOTE (iter-188 plan-phase): body ~115-200 LOC, iter-189+ sub-build.
  % The iter-188 prover does NOT target this body; the iter-188 Lane F
  % target is the consumer `pullback_app_isoTensor_baseMap_isBaseChange`
  % which is closable via `pullback_tildeIso` as a hypothesis-providing
  % helper.
\end{proof}

\begin{lemma}

\leanok
[Pushforward of quasi-coherent modules is quasi-coherent (Stacks 01XJ)]
  \label{lem:pushforward_isQuasicoherent}
  \lean{AlgebraicGeometry.pushforward_isQuasicoherent}
  \textit{Source: Stacks~\href{https://stacks.math.columbia.edu/tag/01XJ}{01XJ}
  (preservation of quasi-coherence under pushforward).}
  Let \(f : X \to Y\) be a quasi-compact, quasi-separated morphism of
  schemes, and let \(\mathcal{F}\) be a quasi-coherent
  \(\mathcal{O}_X\)-module. Then \(f_* \mathcal{F}\) is a quasi-coherent
  \(\mathcal{O}_Y\)-module. Mathlib \texttt{b80f227} ships no
  \texttt{IsQuasicoherent.pushforward} instance at the
  \texttt{Scheme.Modules} layer.
\end{lemma}
\begin{proof}
  Substantive body (Stacks 01XJ): standard adjoint-functor preservation
  under qcqs.

  % NOTE (iter-188 plan-phase): body ~30 LOC, iter-189+ sub-build.
  % Consumer signature `_of_isAffineBase` feeds `(pushforward f).obj F`
  % into `pullback_app_isoTensor`, which now requires
  % `[N.IsQuasicoherent]` — this helper is the qcoh-witness producer.
\end{proof}

\subsection{Iter-189 analogist-licensed unbundle pins}
\label{sec:quot_iter189_unbundle}

The iter-189 \texttt{mathlib-analogist lane-f-isbasechange} verdict
(see \texttt{analogies/lane-f-isbasechange.md}) returned verdict (A)
\emph{STRUCTURAL OK} on the project's use of
\texttt{IsBaseChange.of\_equiv}. The corrective was to \emph{unbundle}
the iter-188 \(\Sigma\)-pair helper's three internal Mathlib gaps into
separately named typed-sorry pins so that each can be targeted
individually (~30-50 LOC apiece) in subsequent iters. Iter-189 Lane F
prover landed the unbundle. The two new pins are recorded below
(parallel to the existing \cref{lem:pullback_tildeIso} pin, which
remains Step 2 of the 3-step chain).

\begin{lemma}

\leanok
[Quasi-coherent module on an affine open is \texttt{tilde} of its sections (Stacks 01I8)]
  \label{lem:tildeIso_of_isQuasicoherent_isAffineOpen}
  \lean{AlgebraicGeometry.tildeIso_of_isQuasicoherent_isAffineOpen}
  \uses{def:quot_canonical_basechange_map}
  \textit{Source: Stacks~\href{https://stacks.math.columbia.edu/tag/01I8}{01I8}
  (a quasi-coherent module on an affine scheme is \(\widetilde M\) for
  some module).}
  Let \(X\) be a scheme, \(\mathcal N\) a quasi-coherent
  \(\mathcal O_X\)-module, and \(V \subset X\) an affine open. Then the
  pullback of \(\mathcal N\) along \texttt{hV.fromSpec} is canonically
  isomorphic to \texttt{tilde} of its sections:
  \[
    (\mathtt{Scheme.Modules.pullback}\,\mathtt{hV.fromSpec}).\mathrm{obj}\,
      \mathcal N
    \;\cong\;
    \widetilde{\,\Gamma(\mathcal N, V)\,}
    \quad\text{(on \(\mathrm{Spec}\,\Gamma(X, V)\)).}
  \]
  Mathlib \texttt{b80f227} ships the analog for \emph{free}
  quasi-coherent modules (\texttt{isIso\_fromTildeΓ\_of\_presentation}
  partial), but the general statement
  ``\([N.IsQuasicoherent] \Rightarrow N|_V \cong \widetilde{\Gamma(N,V)}\)
  on an affine open \(V\)'' (Stacks 01I8 \emph{lemma-affine-qc-module-iff})
  is unowned at the \texttt{Scheme.Modules} layer.
\end{lemma}
\begin{proof}
  \uses{def:quot_canonical_basechange_map}
  Substantive body (Stacks 01I8): the QC sheaf on an affine open is
  determined by its global sections — apply the \texttt{tilde.adjunction}
  counit at the affine open and use that \texttt{tilde} is fully faithful
  on quasi-coherent objects.

  % NOTE (iter-189 plan-phase): body ~20-40 LOC, iter-190+ sub-build.
  % This is Step 1 of the 3-pin unbundle of `_sectionLinearEquiv`. The
  % iter-190 prover does NOT target this body; the iter-190 Lane F prover
  % targets Step 3 (`lem:pullback_of_openImmersion_iso_restrict` below),
  % which is the easiest of the 3 pins per the analogist verdict.
\end{proof}

\begin{lemma}

\leanok
[Section-level transport for pullback along an affine open's
\texttt{fromSpec}]
  \label{lem:pullback_of_openImmersion_iso_restrict}
  \lean{AlgebraicGeometry.pullback_of_openImmersion_iso_restrict}
  \uses{def:quot_canonical_basechange_map}
  \textit{Source: Project-bespoke transport, formed from
  Stacks~\href{https://stacks.math.columbia.edu/tag/01HH}{01HH} +
  \texttt{tilde.isoTop}.}
  Let \(Y\) be a scheme, \(\mathcal N\) an \(\mathcal O_Y\)-module, and
  \(U \subset Y\) an affine open. Equip \(\Gamma((\mathrm{Spec}\,\Gamma(Y, U)),
  \top)\) with the \(\Gamma(Y, U)\)-algebra structure inherited from
  \texttt{Scheme.ΓSpecIso} (the canonical
  \(\Gamma(\mathrm{Spec}\,R, \top) \cong R\) for any commutative ring \(R\)),
  and equip
  \(\Gamma((\text{pullback}\,\mathtt{hU.fromSpec}).\mathrm{obj}\,\mathcal N,
  \top)\) with the \(\Gamma(Y, U)\)-module structure inherited via
  \texttt{Module.compHom}. Then there is a canonical \(\Gamma(Y, U)\)-linear
  isomorphism between this pullback-section module and the original section
  module \(\Gamma(\mathcal N, U)\):
  \[
    \Gamma((\text{pullback}\,\mathtt{hU.fromSpec}).\mathrm{obj}\,\mathcal N,
      \top)
    \;\simeq_{\Gamma(Y, U)}\;
    \Gamma(\mathcal N, U).
  \]
  The map identifies a section of the pulled-back sheaf at \(\top\) of
  \(\mathrm{Spec}\,\Gamma(Y, U)\) with the corresponding section of
  \(\mathcal N\) on the affine open \(U \subset Y\). This pin is the
  iter-190+ prover target as Step 3 of the iter-189 \(\Sigma\)-pair
  unbundle.
\end{lemma}
\begin{proof}
  \uses{def:quot_canonical_basechange_map}
  Substantive body: chain \texttt{hU.toSpecΓ\_fromSpec} (identifies
  \texttt{pullback hU.fromSpec} with \texttt{restrict\_to\_U} up to
  natural iso) with naturality of \texttt{tilde.isoTop} (extracts the
  \(\top\)-section of \texttt{tilde} as the input module), then transport
  the \(\Gamma(Y, U)\)-linear structure through the chain.

  % NOTE (iter-189 plan-phase): body ~30-50 LOC, iter-190+ sub-build.
  % This is Step 3 of the 3-pin unbundle of `_sectionLinearEquiv`, and
  % is the LIGHTEST of the 3 pins per the analogist verdict (the chain
  % `hU.toSpecΓ_fromSpec` + `tilde.isoTop` naturality is a routine but
  % multi-step compositional transport). The iter-190 Lane F prover
  % targets this pin axiom-clean.
\end{proof}

\subsection{Iter-195 Beck-Chevalley 6-stage decomposition + (N1)--(N4) substrate}
\label{sec:quot_iter195_bc_substrate}

The iter-195 plan-phase refactor \texttt{lane-f-step12-sigma-pair}
upgraded the helpers \texttt{step1} and \texttt{step2} of the
Tilde-isoTop route into \(\Sigma\)-pair components carrying
iso-characterizing identities \texttt{\_step1\_apply} and
\texttt{\_step2\_apply}. With these identities in hand, the
Beck-Chevalley intertwining of
\cref{def:pullback_app_isoTensor_sigma} unfolds into the six named
stages spelled out in the prose of that definition; mirror copy at
\texttt{AlgebraicJacobian/Picard/QuotScheme.lean:999-1037}. Stage~1
is closed axiom-clean by the iter-195 Lane F prover; Stages 2--6 are
gated on four substrate naturality lemmas (N1)--(N4) for the
project-side helper \texttt{Scheme.Modules.baseMap}, introduced here
as the iter-196+ prover targets.

The four substrate lemmas all follow a single logical pattern. The
project's \texttt{Scheme.Modules.baseMap} is defined from the unit of
\texttt{pullbackPushforwardAdjunction} (Mathlib
\texttt{Mathlib.AlgebraicGeometry.Sheaves.Modules.Pullback}), so
(N1)--(N3) record naturality squares of that adjunction unit at
triples of morphisms / module homomorphisms / propositional
equalities, and (N4) records the inversion of the \texttt{step3}
component built from \texttt{restrictFunctorIsoPullback} (Mathlib
\texttt{Mathlib.AlgebraicGeometry.OpenImmersion}). Together they
discharge the Beck-Chevalley sequence stage-by-stage. LOC estimate
\(\sim\)20--30 LOC per sub-lemma; total \(\sim\)80--120 LOC of
substrate volume. This is the iter-195 progress-critic-flagged
``volume gap'' relative to the iter-195 plan-phase estimate
(\(\sim\)10--30 LOC for the consumer, which did not budget the
substrate).

\begin{lemma}
[(N1) baseMap naturality in input sheaf]
  \label{lem:baseMap_pullbackComp_apply}
  \lean{AlgebraicGeometry.Scheme.Modules.baseMap_pullbackComp_apply}
  \uses{def:quot_canonical_basechange_map, lem:pullback_tildeIso}
  Let \(g : X \to Y\) be a morphism of schemes, let \(\mathcal N\),
  \(\mathcal N'\) be \(\mathcal O_Y\)-modules, and let
  \(\eta : \mathcal N \to \mathcal N'\) be a morphism of
  \(\mathcal O_Y\)-modules. For affine opens \(U \subseteq Y\) and
  \(V \subseteq X\) with \(g(V) \subseteq U\) and any
  \(x \in \Gamma(\mathcal N, U)\), the section-level component of
  \cref{def:quot_canonical_basechange_map} satisfies
  \[
    \bigl((\text{pullback}\,g).\mathrm{map}\,\eta\bigr).\mathrm{app}\,V\;
    \bigl(\text{baseMap}\,g\,\mathcal N\,e\,x\bigr)
    \;=\;
    \text{baseMap}\,g\,\mathcal N'\,e\,(\eta.\mathrm{app}\,U\,x),
  \]
  i.e.\ \texttt{baseMap} is natural in its module argument. Used in
  Stage~2 of the \cref{def:pullback_app_isoTensor_sigma}
  decomposition to transport \(\text{step1}.\mathrm{inv}\) (a
  morphism of \(\mathcal O_Y\)-modules) across \texttt{baseMap} of
  \(\text{Spec.map}\,\varphi\).
\end{lemma}
\begin{proof}
  \uses{def:quot_canonical_basechange_map}
  Direct from the naturality square of
  \texttt{pullbackPushforwardAdjunction.unit} at the morphism
  \(\eta : \mathcal N \to \mathcal N'\) in the input category, then
  passed to the section level via the definitional unfolding of
  \texttt{Scheme.Modules.baseMap} (which is a section-level reading
  of the unit composed with the canonical pushforward-section map).

  % NOTE (iter-196 plan-phase): body ~20-30 LOC. (N1) of the iter-195
  % 6-stage Beck-Chevalley decomposition; iter-196+ Lane F prover target.
\end{proof}

\begin{lemma}
[(N2) baseMap composition via pullbackComp]
  \label{lem:baseMap_pullback_comp_apply}
  \lean{AlgebraicGeometry.Scheme.Modules.baseMap_pullback_comp_apply}
  \uses{def:quot_canonical_basechange_map, lem:pullback_tildeIso}
  Let \(f : X \to Y\) and \(g : Y \to Z\) be morphisms of schemes and
  let \(\mathcal N\) be an \(\mathcal O_Z\)-module. The
  Mathlib-canonical iso
  \(\text{pullbackComp}\,f\,g :
    (\text{pullback}\,f) \circ (\text{pullback}\,g)
    \cong \text{pullback}\,(f \circ g)\)
  intertwines \texttt{baseMap}: for affine opens \(W \subseteq Z\),
  \(V \subseteq X\) with \(f(g(V)) \subseteq W\) and any
  \(x \in \Gamma(\mathcal N, W)\),
  \[
    ((\text{pullbackComp}\,f\,g)\,\mathcal N).\mathrm{hom}.\mathrm{app}\,V\;
    \bigl(\text{baseMap}\,f\,((\text{pullback}\,g).\mathrm{obj}\,\mathcal N)\,
       e_f\,(\text{baseMap}\,g\,\mathcal N\,e_g\,x)\bigr)
    \;=\;
    \text{baseMap}\,(f \circ g)\,\mathcal N\,e\,x.
  \]
  In words: applying \texttt{baseMap} along \(g\) then along \(f\),
  routed through the canonical composition iso, yields the
  \texttt{baseMap} along the composite \(f \circ g\). Used in Stages~3
  and~5 of the \cref{def:pullback_app_isoTensor_sigma}
  decomposition.
\end{lemma}
\begin{proof}
  \uses{def:quot_canonical_basechange_map}
  The unit of a composition of left adjoints
  \(F \circ G \dashv R \circ R'\) satisfies the canonical composition
  identity \(\eta_{FG} = R'(\eta_F) \circ \eta_G\) (cf.\ Mathlib
  \texttt{Adjunction.comp} unit decomposition). Specialised to the
  pullback-pushforward adjunction at the triple of schemes
  \(X \xrightarrow{f} Y \xrightarrow{g} Z\) and transported to the
  section level via \texttt{Scheme.Modules.baseMap}'s definition,
  this yields the stated formula.

  % NOTE (iter-196 plan-phase): body ~25-35 LOC. (N2) of the iter-195
  % 6-stage Beck-Chevalley decomposition; iter-196+ Lane F prover target.
  % Used twice in the decomposition (Stages 3 and 5).
\end{proof}

\begin{lemma}
[(N3) baseMap transport via pullbackCongr]
  \label{lem:baseMap_pullbackCongr_apply}
  \lean{AlgebraicGeometry.Scheme.Modules.baseMap_pullbackCongr_apply}
  \uses{def:quot_canonical_basechange_map, lem:pullback_tildeIso}
  Let \(f_1, f_2 : X \to Y\) be morphisms of schemes with
  \(f_1 = f_2\) as a propositional equality \(h_{eq}\), and let
  \(\mathcal N\) be an \(\mathcal O_Y\)-module. The Mathlib-canonical
  equality-transport iso
  \(\text{pullbackCongr}\,h_{eq} : \text{pullback}\,f_1
    \cong \text{pullback}\,f_2\)
  intertwines \texttt{baseMap}: for affine opens \(U \subseteq Y\),
  \(V \subseteq X\) with \(f_1(V) \subseteq U\) (equivalently
  \(f_2(V) \subseteq U\) via \(h_{eq}\)) and any
  \(x \in \Gamma(\mathcal N, U)\),
  \[
    ((\text{pullbackCongr}\,h_{eq})\,\mathcal N).\mathrm{inv}.\mathrm{app}\,V\;
    \bigl(\text{baseMap}\,f_2\,\mathcal N\,e_2\,x\bigr)
    \;=\;
    \text{baseMap}\,f_1\,\mathcal N\,e_1\,x.
  \]
  I.e.\ \texttt{baseMap} is invariant under the equality-transport iso
  of the morphism argument. Used in Stage~4 of the
  \cref{def:pullback_app_isoTensor_sigma} decomposition to re-associate
  the composite
  \(\text{Spec.map}\,\varphi \circ \mathtt{\_hV.fromSpec}
    = \mathtt{\_hU.fromSpec} \circ g\) under the Σ-pair packaging.
\end{lemma}
\begin{proof}
  \uses{def:quot_canonical_basechange_map}
  By propositional-equality elimination: write
  \(\text{pullbackCongr}\,h_{eq}\) as
  \(\text{eqToIso}\,(\text{congrArg}\,\text{pullback}\,h_{eq})\) and
  unfold \texttt{eqToIso\_app} to reduce to \texttt{Eq.mpr}; the
  definition of \texttt{baseMap} depends on the morphism argument only
  through \texttt{pullbackPushforwardAdjunction.unit}, which is itself
  definitionally invariant under \texttt{Eq.mpr} along
  \(h_{eq}\).

  % NOTE (iter-196 plan-phase): body ~15-25 LOC; the easiest of the
  % four — almost pure equality transport. (N3) of the iter-195 6-stage
  % Beck-Chevalley decomposition; iter-196+ Lane F prover target.
\end{proof}

\begin{lemma}
[(N4) step3 inversion identity for an open immersion]
  \label{lem:baseMap_inv_step3_open_immersion}
  \lean{AlgebraicGeometry.Scheme.Modules.baseMap_inv_step3_open_immersion}
  \uses{def:quot_canonical_basechange_map,
        lem:pullback_of_openImmersion_iso_restrict}
  Let \(Y\) be a scheme, \(U \subseteq Y\) an affine open with
  canonical open immersion
  \(\mathtt{\_hU.fromSpec} : \text{Spec}\,\Gamma(Y, U) \to Y\), and
  \(\mathcal N\) an \(\mathcal O_Y\)-module. Then the iter-189
  Spec-restriction transport iso \texttt{step3} (constructed in
  \cref{lem:pullback_of_openImmersion_iso_restrict}) is the inverse
  of \texttt{baseMap} along \(\mathtt{\_hU.fromSpec}\) at the
  \(\top\)-section:
  \[
    \text{step3}\bigl(
      \text{baseMap}\,\mathtt{\_hU.fromSpec}\,
        ((\text{pullback}\,g).\mathrm{obj}\,\mathcal N)\,
        \text{le\_top}'\,y
    \bigr) \;=\; y
    \qquad\text{for every } y \in \Gamma(\mathcal N, U).
  \]
  This identifies the inverse of \texttt{step3} with the natural
  \texttt{baseMap} action of the open immersion
  \(\mathtt{\_hU.fromSpec}\). Used in Stage~6 of the
  \cref{def:pullback_app_isoTensor_sigma} decomposition, closing the
  intertwining identity.
\end{lemma}
\begin{proof}
  \uses{def:quot_canonical_basechange_map,
        lem:pullback_of_openImmersion_iso_restrict}
  The transport iso \texttt{step3} is built (in
  \cref{lem:pullback_of_openImmersion_iso_restrict}) from
  \texttt{restrictFunctorIsoPullback} (Mathlib
  \texttt{AlgebraicGeometry.OpenImmersion}), which exhibits, for an
  open immersion \(j\), a natural iso between \(j^*\) and the
  restriction functor \(\text{restrict}\,j\). At the \(\top\)-section
  of \(\text{Spec}\,\Gamma(Y, U)\), the inverse of this natural iso is
  exactly \texttt{baseMap}\,\(\mathtt{\_hU.fromSpec}\), because the
  adjunction unit for an open-immersion pullback acts at \(\top\) as
  the canonical restriction-to-section map.

  % NOTE (iter-196 plan-phase): body ~20-30 LOC. (N4) of the iter-195
  % 6-stage Beck-Chevalley decomposition; iter-196+ Lane F prover target.
\end{proof}

\section{Lean encoding}
\label{sec:quot_lean_encoding}

The Lean target file is the new
\texttt{AlgebraicJacobian/Picard/QuotScheme.lean}. The construction
proceeds in three phases.

\paragraph{Phase 1: Hilbert polynomial (\cref{def:hilbert_polynomial}).}
On a finite-type morphism \(\pi : X \to S\) with \(S\) noetherian and a
coherent \(\mathcal{F}\) on \(X\) with proper support over \(S\), the per-fiber
Hilbert polynomial is encoded as a function \(s \mapsto
\Phi_{\mathcal{F},s} \in \mathbb{Q}[\lambda]\). Mathlib provides
\texttt{CategoryTheory.PolynomialModule} and Euler-characteristic
infrastructure on coherent sheaves on noetherian schemes; the polynomial
property of \(m \mapsto \chi(\mathcal{F} \otimes L^{\otimes m})\) requires
Snapper's Lemma, which is a project-side sub-build on top of the
graded-Euler-characteristic infrastructure.

\paragraph{Phase 2: Grassmannian scheme (\cref{def:grassmannian_scheme},
\cref{thm:grassmannian_representable}).} The Grassmannian
\(\mathrm{Gr}_S(V, d)\) is encoded as the relative gluing of \(\binom{r}{d}\)
affine charts \(U^I\) (each isomorphic to a relative affine space
\(\mathbb{A}^{d(r-d)}_S\)) via the Pl\"ucker cocycle \(\theta_{I,J}\). The
gluing backbone is \texttt{AlgebraicGeometry.Scheme.GlueData}. The
universal quotient sheaf \(\mathcal{U}\) is glued by the transition cocycle
\(g_{I,J} = (X^I_J)^{-1}\), and the determinant line bundle
\(\det \mathcal{U}\) gives the Pl\"ucker closed embedding into
\(\mathbb{P}_S(\bigwedge^d V)\). The representability statement uses
\texttt{CategoryTheory.Functor.RepresentableBy} (as in
\cref{chap:Picard_RelativeSpec}) on the universal pair
\((\mathrm{Gr}_S(V, d), \mathcal{U})\).

\paragraph{Phase 3: Quot scheme (\cref{thm:quot_representable}).}
The Quot functor is constructed as a locally closed subfunctor of the
Grassmannian
\(\mathrm{Grass}(W \otimes_{\mathcal{O}_S} \mathrm{Sym}^r V, \Phi(r))\), via:
\begin{enumerate}
  \item the uniform \(m\)-regularity bound from
    \cref{lem:quot_boundedness} (input: Castelnuovo--Mumford theorem,
    \texttt{AlgebraicGeometry.CastelnuovoMumford} -- project-side sub-build);
  \item the Grassmannian embedding \(\alpha\) from
    \cref{lem:quot_alpha_injective};
  \item the flattening stratification (\cref{thm:flattening_stratification_exists}
    of \cref{chap:Picard_FlatteningStratification}, A.2.a) applied to the
    universal quotient on the Grassmannian, producing the locally closed
    stratum \(T_0^\Phi\);
  \item the valuative-criterion upgrade
    (\cref{lem:quot_valuative_criterion}) to closed embedding, hence
    projective \(S\)-scheme.
\end{enumerate}
The reduction \cref{lem:quot_reduction_to_pi_star_W} from the universal
case to the general case is encoded as a closed embedding of representable
subfunctors via the kernel-vanishing scheme.

\paragraph{Cohomology and base change.}
\cref{thm:flat_base_change_cohomology} is the Mathlib-equivalent of
Stacks tag 02KH. Mathlib carries a partial form
(\texttt{AlgebraicGeometry.flat\_base\_change}) for affine bases; the
full \(R^i f_*\) statement is a project-side bridge if not present at the
pinned commit.

\paragraph{Dependencies summary.}
This file directly consumes
\cref{chap:Picard_RelativeSpec} (A.1.a, on disk) for relative-spectrum
gluing infrastructure and
\cref{thm:flattening_stratification_exists} from
\cref{chap:Picard_FlatteningStratification} (A.2.a, written this iter)
for the flattening stratum. Downstream, \cref{chap:Picard_FGAPicRepresentability}
(A.2.c) consumes \cref{thm:quot_representable} via the special case
\(E = \mathcal{O}_C\) on a smooth proper curve \(C/k\), giving the Hilbert
scheme \(\Hilb_{C/k} = \Quot^{\Phi,\mathcal{O}_C(1)}_{\mathcal{O}_C/C/k}\)
and its open subscheme \(\mathrm{Div}_{C/k}\) of relative effective divisors.

\section{Out of scope}
\label{sec:quot_out_of_scope}

This chapter packages \emph{only} the Quot-scheme existence statement and
its proof skeleton via the Grassmannian + flattening-stratification
route. The following are explicitly out of scope here:
\begin{itemize}
  \item \textbf{Flattening stratification} (A.2.a) -- consumed-by, not
    proved-here. Its blueprint home is the new chapter
    \texttt{Picard\_FlatteningStratification.tex}.
  \item \textbf{Hilbert scheme as a separate chapter}. The Hilbert scheme
    \(\Hilb_{X/S}\) is the special case
    \(\Hilb_{X/S} = \Quot_{\mathcal{O}_X/X/S}\) of the Quot scheme and is
    introduced inline (\cref{def:quot_functor}); no separate chapter.
  \item \textbf{Castelnuovo--Mumford \(m\)-regularity theorem} -- input to
    \cref{lem:quot_boundedness}, not proved here. The proof lives in
    Nitsure \S 2; for the Lean encoding it is a project-side sub-build.
  \item \textbf{Semi-continuity and base change} (Nitsure \S 3) -- the
    flat-base-change theorem of cohomology (\cref{thm:flat_base_change_cohomology})
    is stated here as a lemma, but the upstream Mathlib bridge / project-side
    sub-build is not constructed in this chapter.
  \item \textbf{Snapper's Lemma} -- the polynomial-eventually property of
    Euler characteristics, used in \cref{def:hilbert_polynomial}. Cited
    via Nitsure \S 1 and (textbook reference) Hartshorne III.5.2;
    project-side sub-build for the Lean encoding.
  \item \textbf{Identification with \(\mathbb{P}^n\) and relative
    Grassmannian variants} (Nitsure \S 1, Exercises (1)--(4)) --
    optional applications, not consumed by Route A.
  \item \textbf{Hironaka's counterexample to representability of
    \(\Hilb_{X/\mathbb{C}}\) in the non-projective case} (Nitsure \S 1,
    Exercise (8)) -- pedagogical only; Route A operates entirely in the
    projective setting.
  \item \textbf{Group-scheme structure on \(\Pic_{C/k}\)} (A.2.c) -- the
    abelian-group-scheme refinement is downstream of the Quot existence
    statement; its blueprint home is
    \cref{chap:Picard_FGAPicRepresentability}.
  \item \textbf{Identity-component \(\Pic^0\) and abelian-variety
    structure} (A.3) -- a separate Route A sub-row, fed by but not
    constructed in this chapter.
\end{itemize}
